\documentclass[11pt,a4paper]{article}
\usepackage[utf8]{inputenc}
\usepackage[T1]{fontenc}
\usepackage[english]{babel}
\usepackage{amsmath,amssymb,amsthm}
\usepackage{geometry}
\usepackage{graphicx}
\usepackage{fancyhdr}
\usepackage{titlesec}
\usepackage{enumitem}
\usepackage{array}
\usepackage{booktabs}

\geometry{margin=1in}

\pagestyle{fancy}
\fancyhf{}
\fancyhead[C]{\small The Collected Mathematical Papers of Arthur Cayley, Vol.~XI}
\fancyfoot[C]{\thepage}
\renewcommand{\headrulewidth}{0.4pt}

\titleformat{\section}{\large\bfseries}{\thesection.}{0.5em}{}
\titleformat{\subsection}{\normalsize\bfseries}{\thesubsection.}{0.5em}{}

\theoremstyle{definition}
\newtheorem*{remark}{Remark}

\begin{document}

\begin{center}
{\Large\bfseries The Collected Mathematical Papers of Arthur Cayley}\\[6pt]
{\large Vol.~XI}\\[12pt]
{\large Pages 551--600}\\[24pt]
\end{center}

\tableofcontents
\newpage

% ============================================================
% PAPER 790: GEOMETRY (continued from previous pages)
% ============================================================
\section*{790. GEOMETRY. [Continued]}
\addcontentsline{toc}{section}{790. Geometry (continued)}

\subsection*{\S\,5. Cubical curves (concluded)}
\addcontentsline{toc}{subsection}{\S\,5. Cubical curves (concluded)}

In the last four examples, the curve is one of the cubical curves called the\break divergent parabolas: (iv) is a mere numerical example of (v), and (vi), (vii), (viii) are in Newton's language the parabola cum ovali, punctata, and nodata respectively. When $a$, $b$, $c$ are all equal, or the form is $y^2=(x-c)^3$, we have a cuspidal form, Newton's parabola cuspidata, otherwise the semicubical parabola.

(viii) As an example of a curve given by an implicit equation, suppose the equation is
\[
x^2 + y^2 - 3xy = 0;
\]
this is a nodal cubic curve, the node at the origin, and the axes touching the two branches respectively (fig.~7). An easy mode of tracing it is to express $x$, $y$ each of them in terms of a variable $\theta$,
\[
x = \frac{3\theta}{1+\theta^3},\qquad y = \frac{3\theta^2}{1+\theta^3},
\]
but it is instructive to trace the curve directly from its equation.

\subsection*{\S\,5. Tangent by algebraical process}
\addcontentsline{toc}{subsection}{\S\,5. Tangent by algebraical process}

5. It may be remarked that the purely algebraical process, which is, in fact, that employed in finding a differential coefficient $\dfrac{dy}{dx}$, if applied directly to the equation of the curve, determines the point consecutive to any given point of the curve, that is, the direction of the curve at such given point, or, what is the same thing, the direction of the tangent at that point. In fact, if $\alpha$, $\beta$ are the coordinates of any point on a curve $f(x,y)=0$, then writing in the equation of the curve $x=\alpha+h$, $y=\beta+k$, and in the resulting equation $f(\alpha+h,\beta+k)=0$, developed in powers of $h$ and $k$, omitting the term $f(\alpha,\beta)$, which vanishes, and the terms containing the second and higher powers of $h$, $k$, we have a linear equation $Ah+Bk=0$, which determines the ratio of the increments $h$, $k$. Of course, in the analytical development of the theory, we translate this into the notation of the differential calculus; but the question presents itself, and is thus seen to be solvable, as soon as it is attempted to trace a curve from its equation.

\subsection*{Descriptive and Metrical Geometry}
\addcontentsline{toc}{subsection}{Descriptive and Metrical Geometry}

\textbf{6.} A geometrical proposition is either \textit{descriptive} or \textit{metrical}: in the former case it is altogether independent of the idea of magnitude (length, inclination, \&c.); in the latter case it has reference to this idea. It is to be noticed that, although the method of coordinates seems to be by its inception essentially metrical, and we can hardly, except by metrical considerations, connect an equation with the curve which it represents (for instance, even assuming it to be known that an equation $Ax+By+C=0$ represents a line, yet if it be asked what line, the only form of answer is, that it is the line cutting the axes at distances from the origin $-C\div A$, $-C\div B$ respectively), yet in dealing by this method with descriptive propositions, we are, in fact, eminently free from all metrical considerations.

\textbf{7.} It is worth while to illustrate this by the instance of the well-known theorem of the radical centre of three circles. The theorem is that, given any three circles $A$, $B$, $C$ (fig.~8), the common chords $\alpha\alpha'$, $\beta\beta'$, $\gamma\gamma'$ of the three pairs of circles meet in a point.

The geometrical proof is metrical throughout:

Take the point of intersection of $\alpha\alpha'$, $\beta\beta'$, and joining this with $\gamma'$, suppose that $\gamma'O$ does not pass through $\gamma$, but that it meets the circles $A$, $B$ in two distinct points $\gamma_1$, $\gamma_2$ respectively. We have then the known metrical property of intersecting chords of a circle; viz. in the circle $C$, where $\alpha\alpha'$, $\beta\beta'$ are chords meeting at a point $O$,
\[
O\alpha \cdot O\alpha' = O\beta \cdot O\beta',
\]
where, as well as in what immediately follows $O\alpha$, \&c., denote, of course, lengths or distances.

Similarly in the circle $A$,
\[
O\beta \cdot O\beta' = O\gamma_1 \cdot O\gamma',
\]
and in the circle $B$,
\[
O\alpha \cdot O\alpha' = O\gamma_2 \cdot O\gamma'.
\]
Consequently $O\gamma_1 \cdot O\gamma' = O\gamma_2 \cdot O\gamma'$, that is, $O\gamma_1 = O\gamma_2$, or the points $\gamma_1$ and $\gamma_2$ coincide; that is, they each coincide with $\gamma$.

\newpage
\noindent\textsc{Page~553.} We contrast this with the analytical method.

Here it only requires to be known that an equation $Ax+By+C=0$ represents a line, and an equation $x^2+y^2+Ax+By+C=0$ represents a circle. $A$, $B$, $C$ have, in the two cases respectively, metrical significations but these we are not concerned with. Using $S$ to denote the function $x^2+y^2+Ax+By+C$, the equation of a circle is $S=0$, where $S$ stands for its value; more briefly, we say the equation is $S=x^2+y^2+Ax+By+C=0$. Let the equation of any other circle be $S'=x^2+y^2+A'x+B'y+C'=0$; the equation $S-S'=0$ is a linear equation: $S-S'$ is, in fact, $=(A-A')x+(B-B')y+C-C'$: and it thus represents a line; this equation is satisfied by the coordinates of each of the points of intersection of the two circles (for at each of these points $S=0$ and $S'=0$, therefore also $S-S'=0$); hence the equation $S-S'=0$ is that of the line joining the two points of intersection of the two circles, or say it is the equation of the common chord of the two circles. Considering then a third circle $S''=x^2+y^2+A''x+B''y+C''=0$, the equations of the common chords are $S-S'=0$, $S-S''=0$, $S'-S''=0$ (each of these a linear equation); at the intersection of the first and second of these lines $S=S'$ and $S=S''$, therefore also $S'=S''$, or the equation of the third line is satisfied by the coordinates of the point in question; that is, the three chords intersect in a point $O$, the coordinates of which are determined by the equations $S=S'=S''$.

It further appears that, if the two circles $S=0$, $S'=0$ do not intersect in any real points, they must be regarded as intersecting in two imaginary points, such that the line joining them is the real line represented by the equation $S-S'=0$; or that two circles, whether their intersections be real or imaginary, have always a real common chord (or radical axis), and that for any three circles the common chords intersect in a point (of course real) which is the radical centre. And by this very theorem, given two circles with imaginary intersections, we can, by drawing circles which meet each of them in real points, construct the radical axis of the first-mentioned two circles.

\textbf{8.} The principle employed in showing that the equation of the common chord of two circles is $S-S'=0$ is one of very extensive application, and some more illustrations of it may be given.

Suppose $S=0$, $S'=0$ are lines, that is, let $S$, $S'$ now denote linear functions $Ax+By+C$, $A'x+B'y+C'$, then $S-kS'=0$ ($k$ an arbitrary constant) is the equation of any line passing through the point of intersection of the two given lines. Such a line may be made to pass through any given point, say the point $(x_0,y_0)$; i.e., if $S_0$, $S_0'$ are what $S$, $S'$ respectively become on writing for $(x,y)$ the values $(x_0,y_0)$, then the value of $k$ is $k=S_0\div S_0'$. The equation in fact is $S S_0' - S_0 S' = 0$; and starting from this equation we at once verify it \textit{a posteriori}; the equation is a linear equation satisfied by the values of $(x,y)$ which make $S=0$, $S'=0$; and satisfied also by the values $(x_0,y_0)$; and it is thus the equation of the line in question.

If, as before, $S=0$, $S'=0$ represent circles, then ($k$ being arbitrary) $S-kS'=0$ is the equation of any circle passing through the two points of intersection of the two circles; and to make this pass through a given point $(x_0,y_0)$ we have again $k=S_0\div S_0'$. In the particular case $k=1$, the circle becomes the common chord; more accurately, it becomes the common chord together with the line infinity, but this is a question which is not here gone into.

If $S$ denote the general quadric function,
\[
S = ax^2 + 2hxy + by^2 + 2fy + 2gx + c = (a,b,c,f,g,h\,\raisebox{0.5ex}{\text{$\!(x,y,1)$}})^2,
\]
then the equation $S=0$ represents a conic; assuming this, then, if $S'=0$ represents another conic, the equation $S-kS'=0$ represents any conic through the four points of intersection of the two conics.

Returning to the equation
\[
Ax+By+C=0
\]
of a line, if this pass through two given points $(x_1,y_1)$, $(x_2,y_2)$, then we must have $Ax_1+By_1+C=0$, $Ax_2+By_2+C=0$, equations which determine the ratios $A:B:C$, and it thus appears that the equation of the line through the two given points is
\[
x(y_1-y_2) - y(x_1-x_2) + x_1y_2 - x_2y_1 = 0;
\]
or, what is the same thing,
\[
\begin{vmatrix}
x & y & 1 \\
x_1 & y_1 & 1 \\
x_2 & y_2 & 1
\end{vmatrix} = 0.
\]

\newpage
\noindent\textsc{Page~555.} \textbf{9.} The object still being to illustrate the mode of working with coordinates, we consider the theorem of the polar of a point in regard to a circle. Given a circle and a point (fig.~9), we draw through any two lines meeting the circle in the points $A$, $A'$ and $B$, $B'$ respectively, and then taking $Q$ as the intersection of the lines $AB'$ and $A'B$, the theorem is that the locus of the point $Q$ is a right line depending only upon $O$ and the circle, but independent of the particular lines $OAA'$ and $OBB'$.

Taking $O$ as the origin, and for the axes any two lines through $O$ at right angles to each other, the equation of the circle will be
\[
x^2 + y^2 + 2Ax + 2By + C = 0;
\]
and if the equation of the line $OAA'$ is taken to be $y=mx$, then the points $A$, $A'$ are found as the intersections of the straight line with the circle or to determine $x$ we have
\[
x^2(1+m^2) + 2x(A+Bm) + C = 0.
\]

\noindent\textsc{Page~555 (continued).} If $(x_1,y_1)$ are the coordinates of $A$, and $(x_2,y_2)$ of $A'$, then the roots of this equation are $x_1$, $x_2$, whence easily
\[
\frac{1}{x_1} + \frac{1}{x_2} = \frac{2(A+Bm)}{C}.
\]

And similarly, if the equation of the line $OBB'$ is taken to be $y=m'x$, and the coordinates of $B$, $B'$ to be $(x_3,y_3)$ and $(x_4,y_4)$ respectively, then
\[
\frac{1}{x_3} + \frac{1}{x_4} = \frac{2(A+Bm')}{C}.
\]

We have then
\begin{align*}
x(y_1-y_4) - y(x_1-x_4) + x_1y_4 - x_4y_1 &= 0,\\
x(y_2-y_3) - y(x_2-x_3) + x_2y_3 - x_3y_2 &= 0,
\end{align*}
as the equations of the lines $AB'$ and $A'B$ respectively; for the first of these equations, being satisfied if we write therein $(x_1,y_1)$ or $(x_4,y_4)$ for $(x,y)$, is the equation of the line $AB'$: and similarly the second equation is that of the line $A'B$. Reducing by means of the relations $y_1=mx_1$, $y_2=mx_2$, $y_3=m'x_3$, $y_4=m'x_4$, the two equations become
\begin{align*}
x(mx_4-m'x_1) - y(x_1-x_4) + (m'-m)x_1x_4 &= 0,\\
x(mx_3-m'x_2) - y(x_2-x_3) + (m'-m)x_2x_3 &= 0;
\end{align*}
and if we divide the first of these equations by $mx_1x_4$, and the second by $m'mx_2x_3$, and then add, we obtain
\[
\left(\frac{1}{x_1}+\frac{1}{x_2}\right) - \left(\frac{1}{x_3}+\frac{1}{x_4}\right) y - \left[\left(\frac{1}{x_1}+\frac{1}{x_2}\right) - \left(\frac{1}{x_3}+\frac{1}{x_4}\right)\right] mx = 0,
\]
or, what is the same thing,
\[
\frac{y-m'x}{x_1x_2} - \frac{y-mx}{x_3x_4} = 0,
\]
which by what precedes is the equation of a line through the point $Q$. Substituting herein for $\dfrac{1}{x_1}+\dfrac{1}{x_2}$, $\dfrac{1}{x_3}+\dfrac{1}{x_4}$ their foregoing values, the equation becomes
\[
(A+Bm)(y-m'x) + (A+Bm')(y-mx) + m'-m = 0;
\]
that is,
\[
(m-m')(Ax+By+C) = 0;
\]
or finally it is $Ax+By+C=0$, showing that the point $Q$ lies in a line the position of which is independent of the particular lines $OAA'$, $OBB'$ used in the construction.

It is proper to notice that there is no correspondence to each other of the points $A$, $A'$ and $B$, $B'$; the grouping might as well have been $A$, $A'$ and $B'$, $B$; and it thence appears that the line $Ax+By+C=0$ just obtained is in fact the line joining the point $Q$ with the point $R$ which is the intersection of $AB$ and $A'B'$.

\newpage
\noindent\textsc{Page~556.} \textbf{10.} The equation $Ax+By+C=0$ of a line contains in appearance 3, but really only 2 constants (for one of the constants can be divided out), and a line depends accordingly upon 2 parameters, or can be made to satisfy 2 conditions. Similarly, the equation $(a,b,c,f,g,h\,\raisebox{0.5ex}{\text{$\!(x,y,1)$}})^2=0$ of a conic contains really 5 constants, and the equation $(*\,\raisebox{0.5ex}{\text{$\!(x,y,1)$}})^2=0$ of a cubic contains really 9 constants. It thus appears that a cubic can be made to pass through 9 given points, and that the cubic so passing through 9 given points is completely determined. There is, however, a remarkable exception. Considering two given cubic curves $S=0$, $S'=0$, these intersect in 9 points, and through these 9 points we have the whole series of cubics $S-kS'=0$, where $k$ is an arbitrary constant: $k$ may be determined so that the cubic shall pass through a given tenth point, viz. $k=S_0\div S_0'$, if the coordinates are $(x_0,y_0)$, and $S_0$, $S_0'$ denote the corresponding values of $S$, $S'$. The resulting curve $SS_0'-S'S_0=0$ may be regarded as the cubic determined by the conditions of passing through 8 of the 9 points and through the given point $(x_0,y_0)$; and from the equation it thence appears that the curve passes through the remaining one of the 9 points. In other words, we thus have the theorem, any cubic curve which passes through 8 of the 9 intersections of two given cubic curves passes through the 9th intersection.

The applications of this theorem are very numerous; for instance, we derive from it Pascal's theorem of the inscribed hexagon. Consider a hexagon inscribed in a conic. The three alternate sides constitute a cubic, and the other three alternate sides another cubic. The cubics intersect in 9 points, being the 6 vertices of the hexagon, and the 3 Pascalian points, or intersections of the pairs of opposite sides of the hexagon. Drawing a line through two of the Pascalian points, the conic and this line constitute a cubic passing through 8 of the 9 points of intersection, and it therefore passes through the remaining point of intersection; that is, the third Pascalian point; and since obviously this does not lie on the conic, it must lie on the line; that is, we have the theorem that the three Pascalian points (or points of intersection of the pairs of opposite sides) lie on a line.

\subsection*{Metrical Theory}
\addcontentsline{toc}{subsection}{Metrical Theory}

\textbf{11.} The foundation of the metrical theory consists in the simple theorem that if a finite line $PQ$ (fig.~10) be projected upon any other line $OO'$ by lines perpendicular to $OO'$, then the length of the projection $P'Q'$ is equal to the length of $PQ$ multiplied by the cosine of its inclination to $P'Q'$; or, what is the same thing, that the perpendicular distance $PQ'$ of any two parallel lines is equal to the inclined distance $PQ$ multiplied by the cosine of the inclination. It at once follows that the algebraical sum of the projections of the sides of a closed polygon upon any line is $=0$; or, reversing the signs of certain sides, and considering the polygon as consisting of two broken lines, each extending from the same initial to the same terminal point, the sum of the projections of the lines of the first set upon any line is equal to the sum of the projections of the lines of the second set. Observe that, if any line be perpendicular to the line on which the projection is made, then its projection is $=0$.

Thus, if we have a right-angled triangle $PQR$ (fig.~11), where $QR$, $RP$, $QP$ are $=\xi$, $\eta$, $p$ respectively, and whereof the base-angle is $=\alpha$, then projecting successively on the three sides, we have
\[
\xi = p\cos\alpha,\qquad \eta = p\sin\alpha,\qquad p = \xi\cos\alpha + \eta\sin\alpha
\]
and we thence obtain
\[
p^2 = \xi^2+\eta^2,\qquad \cos^2\alpha+\sin^2\alpha=1.
\]
And again, by projecting on a line $Qx'$, inclined at the angle $\alpha'$ to $QR$, we have
\[
p\cos(\alpha-\alpha') = \xi\cos\alpha' + \eta\sin\alpha'
\]
and by substituting for $\xi$, $\eta$ their foregoing values,
\[
\cos(\alpha-\alpha') = \cos\alpha\cos\alpha' + \sin\alpha\sin\alpha'.
\]
It is to be remarked that, assuming only the theory of similar triangles, we have herein a proof of Euclid, Book~I., Prop.~47; in fact, the same as is given Book~VI., Prop.~31; and also a proof of the trigonometrical formula for $\cos(\alpha-\alpha')$. The formulae for $\cos(\alpha+\alpha')$ and $\sin(\alpha+\alpha')$ could be obtained in the same manner.

Draw $PT$ at right angles to $Qx'$, and suppose $QT$, $TP=\xi_1$, $\eta_1$ respectively, so that we have now the quadrilateral $QRPTQ$, or, what is the same thing, the two broken lines $QRP$ and $QTP$, each extending from $Q$ to $P$. Projecting on the four sides successively, we have
\begin{align*}
\xi &= \xi_1\cos\alpha' + \eta_1\sin\alpha',\\
\eta &= \xi_1\sin\alpha' + \eta_1\cos\alpha',\\
\xi_1 &= \xi\cos\alpha' + \eta\sin\alpha',\\
\eta_1 &= -\xi\sin\alpha' + \eta\cos\alpha',
\end{align*}
where the third equation is that previously written $p\cos(\alpha-\alpha') = \xi\cos\alpha' + \eta\sin\alpha$.

\newpage
\noindent\textsc{Page~558.} \subsection*{Equations of Right Line and Circle. Transformation of Coordinates}
\addcontentsline{toc}{subsection}{Equations of Right Line and Circle. Transformation of Coordinates}

\textbf{12.} The required formulae are really contained in the foregoing results. For, in fig.~11, supposing that the axis of $x$ is parallel to $QR$, and taking $a$, $b$ for the coordinates of $Q$, and $(x,y)$ for those of $P$, then we have $\xi$, $\eta = x-a$, $y-b$ respectively and therefore
\begin{gather*}
x-a = p\cos\alpha,\qquad y-b = p\sin\alpha,\\
p^2 = (x-a)^2+(y-b)^2.
\end{gather*}
Writing the first two of these in the form
\[
\frac{x-a}{\cos\alpha} = \frac{y-b}{\sin\alpha} (= p),
\]
we may regard $Q$ as a fixed point, but $P$ as a point moving in the direction $Q$ to $P$, so that $\alpha$ remains constant, and then, omitting the equation $(=p)$, we have a relation between the coordinates $x$, $y$ of the point thus moving in a right line, that is, we have the equation of the line through the given point $(a,b)$ at a given inclination $\alpha$ to the axis of $x$. And, moreover, if, using this equation $(=p)$, we write $x=a+p\cos\alpha$, $y=b+p\sin\alpha$, then we have expressions for the coordinates $x$, $y$ of a point of this line, in terms of the variable parameter $p$.

Again, take the point $T$ to be fixed, but consider the point $P$ as moving in the line $TP$ at right angles to $QT$. If instead of $\xi_1$ we take $p$ for the distance $QT$, then the equation $\xi_1=\xi\cos\alpha'+\eta\sin\alpha'$ will be
\[
(x-a)\cos\alpha' + (y-b)\sin\alpha' = p;
\]
that is, this will be the equation of a line such that its perpendicular distance from the point $(a,b)$ is $=p$, and that the inclination of this distance to the axis of $x$ is $=\alpha'$.

From either form it appears that the equation of a line is, in fact, a linear equation of the form $Ax+By+C=0$. It is important to notice that, starting from this equation, we can determine conversely the $\alpha$ but not the $(a,b)$ of the form of equation which contains these quantities; and in like manner the $\alpha'$ but not the $(a,b)$ or $p$ of the other form of equation. The reason is obvious. In each case $(a,b)$ denote the coordinates of a point, fixed indeed, but which is in the first form any point of the line, and in the second form any point whatever. Thus, in the second form the point from which the perpendicular is let fall may be the origin. Here $(a,b)=(0,0)$, and the equation is $x\cos\alpha'+y\sin\alpha'-p=0$. Comparing this with $Ax+By+C=0$, we have the values of $\cos\alpha'$, $\sin\alpha'$, and $p$.

\textbf{13.} The equation
\[
p^2 = (x-a)^2+(y-b)^2
\]
is an expression for the squared distance of the two points $(a,b)$ and $(x,y)$. Taking as before the point $Q$, coordinates $(a,b)$, as a fixed point, and writing $c$ in the place of $p$, the equation
\[
(x-a)^2+(y-b)^2=c^2
\]
expresses that the point $(x,y)$ is always at a given distance $c$ from the given point $(a,b)$; viz.\ this is the equation of a circle, having $(a,b)$ for the coordinates of its centre, and $c$ for its radius.

The equation is of the form
\[
x^2+y^2+2Ax+2By+C=0,
\]
and here, the number of constants being the same, we can identify the two equations; we find $a=-A$, $b=-B$, $c^2=A^2+B^2-C$, or the last equation is that of a circle having $-A$, $-B$ for the coordinates of its centre, and $\sqrt{A^2+B^2-C}$ for its radius.

\newpage
\noindent\textsc{Page~559.} \textbf{14.} Drawing (fig.~11) $Qy_1$ at right angles at $Qx_1$, and taking $Qx_1$, $Qy_1$ as a new set of rectangular axes, if instead of $\xi_1$, $\eta_1$ we write $x_1$, $y_1$, we have $x_1$, $y_1$ as the new coordinates of the point $P$; and writing also $\alpha$ in place of $\alpha'$, $\alpha$ now denoting the inclination of the axes $Qx_1$ and $Ox$, we have the formulas for transformation between two sets of rectangular axes. These are
\begin{align*}
x-a &= x_1\cos\alpha - y_1\sin\alpha, & x_1 &= (x-a)\cos\alpha + (y-b)\sin\alpha,\\
y-b &= x_1\sin\alpha + y_1\cos\alpha, & y_1 &= -(x-a)\sin\alpha + (y-b)\cos\alpha,
\end{align*}
each set being obviously at once deducible from the other one. In these formulae $(a,b)$ are the $xy$-coordinates of the new origin $Q$, and $\alpha$ is the inclination of $Qx_1$ to $Ox$. It is to be noticed that $Qx_1$, $Qy_1$ are so placed that, by moving to $Q$, and then turning the axes $Ox_1$, $Oy_1$ round $Q$ (through an angle $\alpha$ measured in the sense $Ox$ to $Oy$), the original axes $Ox$, $Oy$ will come to coincide with $Qx_1$, $Qy_1$ respectively. This could not have been done if $Qy_1$ had been drawn (at right angles always to $Qx_1$) in the reverse direction: we should then have had in the formulae $-y_1$ instead of $+y_1$. The new formulae which would be thus obtained are of an essentially distinct form: the analytical test is that in the formulae as written down we can, by giving to $\alpha$ a proper value (in fact, $\alpha=0$), make the $(x-a)$ and $(y-b)$ equal to $x_1$ and $y_1$ respectively; in the other system we could only make them equal to $x_1$, $-y_1$, or $-x_1$, $y_1$ respectively. But for the very reason that the second system can be so easily derived from the first, it is proper to attend exclusively to the first system, that is, always to take the new axes so that the two sets admit of being brought into coincidence.

In the foregoing system of two pairs of equations, the first pair give the original coordinates $x$, $y$ in terms of the new coordinates $x_1$, $y_1$; the second pair the new coordinates $x_1$, $y_1$ in terms of the original coordinates $x$, $y$. The formulae involve $(a,b)$, the original coordinates of the new origin; it would be easy, instead of these, to introduce $(a_1,b_1)$, the new coordinates of the origin. Writing $(a,b)=(0,0)$, we have, of course, the formulae for transformation between two sets of rectangular axes having the same origin, and it is as well to write the formulae in this more simple form; the subsequent transformation to a new origin, but with axes parallel to the original axes, can then be effected without any difficulty.

\newpage
\noindent\textsc{Page~560.} \textbf{15.} All questions in regard to the line may be solved by means of one or other of the foregoing forms
\begin{gather*}
Ax+By+C=0,\\
y=Ax+B,\\
\frac{x-a}{\cos\alpha} = \frac{y-b}{\sin\alpha},\\
(x-a)\cos\alpha' + (y-b)\sin\alpha' - p = 0;
\end{gather*}
or it may be by a comparison of these different forms: thus, using the first form, it has been already shown that the equation of the line through two given points $(x_1,y_1)$, $(x_2,y_2)$ is
\[
x(y_1-y_2)-y(x_1-x_2)+x_1y_2-x_2y_1=0,
\]
or, as this may be written,
\[
\begin{vmatrix} x & y & 1 \\ x_1 & y_1 & 1 \\ x_2 & y_2 & 1 \end{vmatrix} = 0;
\]
a particular case is the equation
\[
\frac{x}{a}+\frac{y}{b}=1
\]
representing the line through the points $(a,0)$ and $(0,b)$, or, what is the same thing, the line meeting the axes of $x$ and $y$ at the distances from the origin $a$ and $b$ respectively. It may be noticed that, in the form $Ax+By+C=0$, $-B\div A$ denotes the tangent of the inclination to the axis of $x$, or we may say that $B\div\sqrt{A^2+B^2}$ and $-A\div\sqrt{A^2+B^2}$ denote respectively the cosine and the sine of the inclination to the axis of $x$. A better form is this: $A\div\sqrt{A^2+B^2}$ and $B\div\sqrt{A^2+B^2}$ denote respectively the cosine and the sine of inclination to the axis of $x$ of the perpendicular upon the line. So, of course, in regard to the form $y=Ax+B$, $A$ is here the tangent of the inclination to the axis of $x$; $1\div\sqrt{A^2+1}$ and $A\div\sqrt{A^2+1}$ are the cosine and sine of this inclination, \&c. It thus appears that the condition, in order that the lines $Ax+By+C=0$ and $A'x+B'y+C'=0$ may meet at right angles, is $AA'+BB'=0$; so when the equations are $y=Ax+B$, $y=A'x+B'$, the condition is $AA'+1=0$, or say the value of $A'$ is
\[
A' = -\frac{1}{A}.
\]

The perpendicular distance of the point $(a,b)$ from the line $Ax+By+C=0$ is $(Aa+Bb+C)\div\sqrt{A^2+B^2}$. In all the formulae involving $\sqrt{A^2+B^2}$ or $\sqrt{A^2+1}$, the radical should be written with the sign $\pm$, which is essentially indeterminate: the like indeterminateness of sign presents itself in the expression for the distance of two points $p=\pm\sqrt{(x-a)^2+(y-b)^2}$; if, as before, the points are $Q$, $P$, and the indefinite line through these is $z'QPz$, then it is the same thing whether we measure off from $Q$ along this line, considered as drawn from $z'$ towards $z$, a positive distance $k$, or along the line considered as drawn reversely from $z$ towards $z'$, the equal negative distance $-k$, and the expression for the distance $p$ is thus properly of the form $\pm k$.

It is interesting to compare expressions which do not involve a radical: thus, in seeking for the expression for the perpendicular distance of the point $(a,b)$ from a given line, let the equation of the given line be taken in the form,
\[
x\cos\alpha + y\sin\alpha - p = 0,
\]
$p$ being the perpendicular distance from the origin, $\alpha$ its inclination to the axis of $x$: the equation of the line may also be written $(x-a)\cos\alpha+(y-b)\sin\alpha-p_1=0$, and we have thence $p_1=p-a\cos\alpha-b\sin\alpha$, the required expression for the distance $p_1$: it is here assumed that $p_1$ is drawn from $(a,b)$ in the same sense as $p$ is drawn from the origin, and the indeterminateness of sign is thus removed.

\newpage
\noindent\textsc{Page~561.} \textbf{16.} As an instance of the mode of using the formulae, take the problem of finding the locus of a point such that its distance from a given point is in a given ratio to its distance from a given line.

We take $(a,b)$ as the coordinates of the given point, and it is convenient to take $(x,y)$ as the coordinates of the variable point, the locus of which is required: it thus becomes necessary to use other letters, say $(X,Y)$, for current coordinates in the equation of the given line. Suppose this is a line such that its perpendicular distance from the origin is $=p$, and that the inclination of $p$ to the axis of $x$ is $=\alpha$; the equation is $X\cos\alpha+Y\sin\alpha-p=0$. In the result obtained in \S~15, writing $(x,y)$ in place of $(a,b)$, it appears that the perpendicular distance of this line from the point $(x,y)$ is
\[
=p-x\cos\alpha-y\sin\alpha;
\]
hence the equation of the locus is
\[
\sqrt{(x-a)^2+(y-b)^2} = e(p-x\cos\alpha-y\sin\alpha),
\]
or say
\[
(x-a)^2+(y-b)^2-e^2(x\cos\alpha+y\sin\alpha-p)^2=0,
\]
an equation of the second order.

\subsection*{The Conics (Parabola, Ellipse, Hyperbola)}
\addcontentsline{toc}{subsection}{The Conics (Parabola, Ellipse, Hyperbola)}

\textbf{17.} The conics or, as they were called, conic sections were originally defined as the sections of a right circular cone; but Apollonius substituted a definition, which is, in fact, that of the last example: the curve is the locus of a point such that its distance from a given point (called the \textit{focus}) is in a given ratio to its distance from a given line (called the \textit{directrix}); and taking the ratio as $e:1$, then $e$ is called the \textit{eccentricity}.

Take $FD$ for the perpendicular from the focus $F$ upon the directrix, and the given ratio being that of $e:1$ ($e>$, $=$, or $<1$, but positive), and let the distance $FD$ be divided at $O$ in the given ratio, say we have $OD=m$, $OF=em$, where $m$ is positive; then the origin may be taken at $O$, the axis $Ox$ being in the direction $OF$ (that is, from $O$ to $F$), and the axis $Oy$ at right angles to it. The distance of the point $(x,y)$ from $F$ is $=\sqrt{(x-em)^2+y^2}$, its distance from the directrix is $=x+m$; the equation therefore is
\[
(x-em)^2+y^2=e^2(x+m)^2;
\]
or, what is the same thing, it is
\[
(1-e^2)x^2-2me(1+e)x+y^2=0.
\]
If $e^2=1$, or, since $e$ is taken to be positive, if $e=1$, this is
\[
y^2-4mx=0,
\]
which is the parabola.

\newpage
\noindent\textsc{Page~562.} If $e^2$ not $=1$, then the equation may be written
\[
(1-e^2)\left\{x-\frac{me}{1-e}\right\}^2+y^2=m^2e^2(1-e^2).
\]
Supposing $e$ positive and $<1$, then, writing $a=\dfrac{me}{1-e^2}$, the equation becomes
\[
(1-e^2)(x-a)^2+y^2=a^2(1-e^2)^2,
\]
that is,
\[
\frac{(x-a)^2}{a^2}+\frac{y^2}{a^2(1-e^2)}=1;
\]
or, changing the origin and writing $b^2=a^2(1-e^2)$, this is
\[
\frac{x^2}{a^2}+\frac{y^2}{b^2}=1,
\]
which is the ellipse.

And similarly if $e$ be positive and $>1$, then writing $a=\dfrac{me}{e^2-1}$, the equation becomes
\[
(1-e^2)(x+a)^2+y^2=a^2(1-e^2)^2,
\]
that is,
\[
\frac{(x+a)^2}{a^2}+\frac{y^2}{a^2(1-e^2)}=1;
\]
or changing the origin and writing $b^2=a^2(e^2-1)$, this is
\[
\frac{x^2}{a^2}-\frac{y^2}{b^2}=1,
\]
which is the hyperbola.

\textbf{18.} The general equation $ax^2+2hxy+by^2+2fy+2gx+c=0$, or as it is written $(a,b,c,f,g,h\,\raisebox{0.5ex}{\text{$\!(x,y,1)$}})^2=0$, may be such that the quadric function breaks up into factors, $=(\alpha x+\beta y+\gamma)(\alpha' x+\beta' y+\gamma')$; and in this case the equation represents a pair of lines, or (it may be) two coincident lines. When it does not so break up, the function can be put in the form $\lambda\{(x-a')^2+(y-b')^2-e^2(x\cos\alpha+y\sin\alpha-p)^2\}$, or, equating the two expressions, there will be six equations for the determination of $\lambda$, $a'$, $b'$, $e$, $p$, $\alpha$; and by what precedes, if $a$, $b'$, $e$, $p$, $\alpha$ are real, the curve is either a parabola, ellipse, or hyperbola. The original coefficients $(a,b,c,f,g,h)$ may be such as not to give any system of real values for $a$, $b'$, $e$, $p$, $\alpha$; but when this is so the equation $(a,b,c,f,g,h\,\raisebox{0.5ex}{\text{$\!(x,y,1)$}})^2=0$ does not represent a real curve\footnote{It is proper to remark that, when $(a,b,c,f,g,h)(x,y,1)^2=0$ does represent a real curve, there are, in fact, four systems of values of $a'$, $b'$, $e$, $p$, $\alpha$, two real, the other two imaginary; we have thus two real equations and two imaginary equations, each of them of the form $(x-a')^2+(y-b')^2=e^2(x\cos\alpha+y\cos\beta-p)^2$, representing each of them one and the same real curve. This is consistent with the assertion of the text that the real curve is in every case represented by a real equation of this form.}; the imaginary curve which it represents is, however, regarded as a conic. Disregarding the special cases of the pair of lines and the twice repeated line, it thus appears that the only real curves represented by the general equation $(a,b,c,f,g,h\,\raisebox{0.5ex}{\text{$\!(x,y,1)$}})^2=0$ are the parabola, the ellipse, and the hyperbola. The circle is considered as a particular case of the ellipse.

The same result is obtained by transforming the equation $(a,b,c,f,g,h)(x,y,1)^2=0$ to new axes. If in the first place the origin be unaltered, then the directions of the new (rectangular) axes $Ox_1$, $Oy_1$ can be found so that $h_1$ (the coefficient of the term $x_1y_1$) shall be $=0$; when this is done, then either one of the coefficients of $x_1^2$, $y_1^2$ is $=0$, and the curve is then a parabola, or neither of these coefficients is $=0$, and the curve is then an ellipse or hyperbola, according as the two coefficients are of the same sign or of opposite signs.

\newpage
\noindent\textsc{Page~563.} \textbf{19.} The curves can be at once traced from their equations:
\begin{align*}
y^2 &= 4mx, && \text{for the parabola (fig.~13),}\\
\frac{x^2}{a^2}+\frac{y^2}{b^2} &= 1, && \text{for the ellipse (fig.~14),}\\
\frac{x^2}{a^2}-\frac{y^2}{b^2} &= 1, && \text{for the hyperbola (fig.~15);}
\end{align*}
and it will be noticed how the form of the last equation puts in evidence the two asymptotes
\[
\frac{x}{a}\pm\frac{y}{b}=0
\]
of the hyperbola. Referred to the asymptotes (as a set of oblique axes) the equation of the hyperbola takes the form $xy=c$; and in particular, if in this equation the axes are at right angles, then the equation represents the rectangular hyperbola referred to its asymptotes as axes.

\subsection*{Tangent, Normal, Circle and Radius of Curvature, \&c.}
\addcontentsline{toc}{subsection}{Tangent, Normal, Circle and Radius of Curvature, \&c.}

\textbf{20.} There is great convenience in using the language and notation of the infinitesimal analysis; thus we consider on a curve a point with coordinates $(x,y)$, and a consecutive point the coordinates of which are $(x+dx, y+dy)$, or again a second consecutive point with coordinates $(x+dx+\tfrac12d^2x, y+dy+\tfrac12d^2y)$, \&c.; and in the final results the ratios of the infinitesimals must be replaced by differential coefficients in the proper manner; thus, if $x$, $y$ are considered as given functions of a parameter $\theta$, then $dx$, $dy$ have in fact the values $\dfrac{dx}{d\theta}d\theta$, $\dfrac{dy}{d\theta}d\theta$, and (only the ratio being really material) they may in the result be replaced by $\dfrac{dx}{d\theta}$, $\dfrac{dy}{d\theta}$. This includes the case where the equation of the curve is given in the form $y=\phi(x)$; $\theta$ is here $=x$, and the increments $dx$, $dy$ are in the result to be replaced by $1$, $\dfrac{dy}{dx}$. So also with the infinitesimals of the higher orders $d^2x$, \&c.

\textbf{21.} The tangent at the point $(x,y)$ is the line through this point and the consecutive point $(x+dx, y+dy)$; hence, taking $\xi$, $\eta$ as current coordinates, the equation is
\[
\frac{\xi-x}{dx}=\frac{\eta-y}{dy},
\]
an equation which is satisfied on writing therein $\xi$, $\eta=(x,y)$ or $=(x+dx, y+dy)$. The equation may be written
\[
\eta-y=\frac{dy}{dx}(\xi-x),
\]
$\dfrac{dy}{dx}$ being now the differential coefficient of $y$ in regard to $x$; and this form is applicable whether $y$ is given directly as a function of $x$, or in whatever way $y$ is in effect given as a function of $x$: if as before $x$, $y$ are given each of them as a function of $\theta$, then the value of $\dfrac{dy}{dx}$ is $=\left(\dfrac{dy}{d\theta}\right)\div\left(\dfrac{dx}{d\theta}\right)$, which is the result obtained from the original form on writing therein $\dfrac{dx}{d\theta}$, $\dfrac{dy}{d\theta}$, for $dx$, $dy$ respectively.

So again, when the curve is given by an equation $u=0$ between the coordinates $(x,y)$, then $\dfrac{dy}{dx}$ is obtained from the equation $\dfrac{du}{dx}+\dfrac{du}{dy}\dfrac{dy}{dx}=0$. But here it is more elegant, using the original form, to eliminate $dx$, $dy$ by the formula $\dfrac{du}{dx}dx+\dfrac{du}{dy}dy$; we thus obtain the equation of the tangent in the form
\[
(\xi-x)\frac{du}{dx}+(\eta-y)\frac{du}{dy}=0.
\]
For example, in the case of the ellipse $\dfrac{x^2}{a^2}+\dfrac{y^2}{b^2}=1$, the equation is
\[
\frac{x}{a^2}(\xi-x)+\frac{y}{b^2}(\eta-y)=0;
\]
or reducing by means of the equation of the curve the equation of the tangent is
\[
\frac{x\xi}{a^2}+\frac{y\eta}{b^2}=1.
\]
The normal is a line through the point at right angles to the tangent; the equation therefore is
\[
(\xi-x)\,dx+(\eta-y)\,dy=0,
\]
where $dx$, $dy$ are to be replaced by their proportional values as before.

\newpage
\noindent\textsc{Page~565.} \textbf{22.} The circle of curvature is the circle through the point and two consecutive points of the curve. Taking the equation to be
\[
(\xi-\alpha)^2+(\eta-\beta)^2=\rho^2,
\]
the values of $\alpha$, $\beta$ are given by
\[
\alpha=x-\frac{dy(d\!x^2+dy^2)}{dx\,d^2y-dy\,d^2x},\qquad
\beta=y+\frac{dx(d\!x^2+dy^2)}{dx\,d^2y-dy\,d^2x},
\]
and we then have
\[
\rho^2=(x-\alpha)^2+(y-\beta)^2=\frac{(d\!x^2+dy^2)^3}{(dx\,d^2y-dy\,d^2x)^2}.
\]
In the case where $y$ is given directly as a function of $x$, then, writing for shortness $p=\dfrac{dy}{dx}$, $q=\dfrac{d^2y}{dx^2}$, this is $\rho^2=\dfrac{(1+p^2)^3}{q^2}$, or, as the equation is usually written, $\rho=\dfrac{(1+p^2)^{\frac32}}{q}$, the radius of curvature, considered to be positive or negative according as the curve is concave or convex to the axis of $x$.

It may be added that the centre of curvature is the intersection of the normal by the consecutive normal.

The locus of the centre of curvature is the \textit{evolute}. If from the expressions of $\alpha$, $\beta$ regarded as functions of $x$ we eliminate $x$, we have thus an equation between $(\alpha,\beta)$, which is the equation of the evolute.

\subsection*{Polar Coordinates}
\addcontentsline{toc}{subsection}{Polar Coordinates}

\textbf{23.} The position of a point may be determined by means of its distance from a fixed point and the inclination of this distance to a fixed line through the fixed point. Say we have $r$ the distance from the origin, and $\theta$ the inclination of $r$ to the axis of $x$; $r$ and $\theta$ are then the polar coordinates of the point, $r$ the radius vector, and $\theta$ the inclination. These are immediately connected with the Cartesian coordinates $x$, $y$ by the formulae $x=r\cos\theta$, $y=r\sin\theta$; and the transition from either set of coordinates to the other can thus be made without difficulty. But the use of polar coordinates is very convenient,---as well in reference to certain classes of questions relating to curves of any kind (for instance, in the dynamics of central forces) as in relation to curves having in regard to the origin the symmetry of the regular polygon (curves such as that represented by the equation $r=\cos m\theta$), and also in regard to the class of curves called spirals, where the radius vector $r$ is given as an algebraical or exponential function of the inclination $\theta$.

\subsection*{Trilinear Coordinates}
\addcontentsline{toc}{subsection}{Trilinear Coordinates}

\textbf{24.} Consider a fixed triangle $ABC$, and (regarding the sides as indefinite lines) suppose for a moment that $p$, $q$, $r$ denote the distances of a point $P$ from the sides $BC$, $CA$, $AB$ respectively, these distances being measured either perpendicularly to the several sides, or each of them in a given direction. To fix the ideas each distance may be considered as positive for a point inside the triangle, and the sign is thus fixed for any point whatever. There is then an identical relation between $p$, $q$, $r$: if $a$, $b$, $c$ are the lengths of the sides, and the distances are measured perpendicularly thereto, the relation is $ap+bq+cr=\text{twice the area of triangle}$. But taking $x$, $y$, $z$ proportional to $p$, $q$, $r$, or if we please proportional to given multiples of $p$, $q$, $r$, then only the ratios of $x$, $y$, $z$ are determined; their absolute values remain arbitrary. But the ratios of $p$, $q$, $r$, and consequently also the ratios of $x$, $y$, $z$ determine, and that uniquely, the point; and it being understood that only the ratios are attended to, we say that $(x,y,z)$ are the coordinates of the point. The equation of a line has thus the form $ax+by+cz=0$, and generally that of a curve of the $n$th order is a homogeneous equation of this order between the coordinates, $(*\,\raisebox{0.5ex}{\text{$\!(x,y,z)$}})^n=0$. The advantage over Cartesian coordinates is in the greater symmetry of the analytical forms, and in the more convenient treatment of the line infinity and of points at infinity. The method includes that of Cartesian coordinates, the homogeneous equation in $x$, $y$, $z$ is, in fact, an equation in $\dfrac{x}{z}$, $\dfrac{y}{z}$ which two quantities may be regarded as denoting Cartesian coordinates; or, what is the same thing, we may in the equation write $z=1$. It may be added that, if the trilinear coordinates $(x,y,z)$ are regarded as the Cartesian coordinates of a point of space, then the equation is that of a cone having the origin for its vertex; and conversely that such equation of a cone may be regarded as the equation in trilinear coordinates of a plane curve.

\newpage
\noindent\textsc{Page~567.} \subsection*{General Point-Coordinates. Line-Coordinates}
\addcontentsline{toc}{subsection}{General Point-Coordinates. Line-Coordinates}

\textbf{25.} All the coordinates considered thus far are point-coordinates. More generally, any two quantities (or the ratios of three quantities) serving to determine the position of a point in the plane may be regarded as the coordinates of the point; or, if instead of a single point they determine a system of two or more points, then as the coordinates of the system of points. But, as noticed under Curve, [785], there are also line-coordinates serving to determine the position of a line; the ordinary case is when the line is determined by means of the ratios of three quantities $\xi$, $\eta$, $\zeta$ (correlative to the trilinear coordinates $x$, $y$, $z$). A linear equation $a\xi+b\eta+c\zeta=0$ represents then the system of lines such that the coordinates of each of them satisfy this relation, in fact, all the lines which pass through a given point; and it is thus regarded as the line-equation of this point; and generally a homogeneous equation $(*\,\raisebox{0.5ex}{\text{$\!(\xi,\eta,\zeta)$}})^n=0$ represents the curve which is the envelope of all the lines the coordinates of which satisfy this equation, and it is thus regarded as the line-equation of this curve.

\section*{II. Solid Analytical Geometry (\S\S\ 26--40)}
\addcontentsline{toc}{section}{II. Solid Analytical Geometry}

\textbf{26.} We are here concerned with points in space, the position of a point being determined by its three coordinates $x$, $y$, $z$. We consider three coordinate planes, at right angles to each other, dividing the whole of space into eight portions called octants, the coordinates of a point being the perpendicular distances of the point from the three planes respectively, each distance being considered as positive or negative according as it lies on the one or the other side of the plane. Thus the coordinates in the eight octants have respectively the signs
\begin{center}
\begin{tabular}{c|cccccccc}
& I & II & III & IV & V & VI & VII & VIII \\
\hline
$x$ & $+$ & $-$ & $-$ & $+$ & $+$ & $-$ & $-$ & $+$ \\
$y$ & $+$ & $+$ & $-$ & $-$ & $+$ & $+$ & $-$ & $-$ \\
$z$ & $+$ & $+$ & $+$ & $+$ & $-$ & $-$ & $-$ & $-$ \\
\end{tabular}
\end{center}

The positive parts of the axes are usually drawn as in fig.~16, which represents a point $P$, the coordinates of which have the positive values $OM$, $MN$, $NP$.

\newpage
\noindent\textsc{Page~568.} \textbf{27.} It may be remarked, as regards the delineation of such solid figures, that if we have in space three lines at right angles to each other, say $Oa$, $Ob$, $Oc$, of equal lengths, then it is possible to project these by parallel lines upon a plane in such wise that the projections $Oa'$, $Ob'$, $Oc'$ shall be at given inclinations to each other, and that these lengths shall be to each other in given ratios: in particular, the two lines $Oa'$, $Oc'$ may be at right angles to each other, and their lengths equal, the direction of $Ob'$, and its proportion to the two equal lengths $Oa'$, $Oc'$, being arbitrary. It thus appears that we may as in the figure draw $Ox$, $Oz$ at right angles to each other, and $Oy$ in an arbitrary direction; and moreover represent the coordinates $x$, $z$ on equal scales, and the remaining coordinate $y$ on an arbitrary scale (which may be that of the other two coordinates $x$, $z$, but is in practice usually smaller). The advantage, of course, is that a figure in one of the coordinate planes $xz$ is represented in its proper form without distortion; but it may be in some cases preferable to employ the isometrical projection, wherein the three axes are represented by lines inclined to each other at angles of $120^\circ$, and the scales for the coordinates are equal (fig.~17).

For the delineation of a surface of a tolerably simple form, it is frequently sufficient to draw (according to the foregoing projection) the sections by the coordinate planes; and in particular, when the surface is symmetrical in regard to the coordinate planes, it is sufficient to draw the quarter-sections belonging to a single octant of the surface; thus fig.~18 is a convenient representation of an octant of the wave surface. Or a surface may be delineated by means of a series of parallel sections, or (taking these to be the sections by a series of horizontal planes) say by a series of contour lines. Of course, other sections may be drawn or indicated, if necessary. For the delineation of a curve, a convenient method is to represent, as above, a series of the points thereof, each point being accompanied by the ordinate $PN$, which serves to refer the point to the plane of $xy$; this is in effect a representation of each point $P$ of the curve, by means of two points $P$, $N$ such that the line $PN$ has a fixed direction. Both as regards curves and surfaces, the employment of stereographic representations is very interesting.

\textbf{28.} In plane geometry, reckoning the line as a curve of the first order, we have only the point and the curve. In solid geometry, reckoning a line as a curve of the first order, and the plane as a surface of the first order, we have the point, the curve, and the surface; but the increase of complexity is far greater than would hence at first sight appear. In plane geometry a curve is considered in connexion with lines (its tangents); but in solid geometry the curve is considered in connexion with lines and planes (its tangents and osculating planes), and the surface also in connexion with lines and planes (its tangent lines and tangent planes); there are surfaces arising out of the line---cones, skew surfaces, developables, doubly and triply infinite systems of lines, and whole classes of theories which have nothing analogous to them in plane geometry: it is thus a very small part indeed of the subject which can be even referred to in the present article.

In the case of a surface, we have between the coordinates $(x,y,z)$ a single, or say a onefold relation, which can be represented by a single relation $f(x,y,z)=0$; or we may consider the coordinates expressed each of them as a given function of two variable parameters $p$, $q$; the form $z=f(x,y)$ is a particular case of each of these modes of representation; in other words, we have in the first mode $f(x,y,z)=z-f(x,y)$, and in the second mode $x=p$, $y=q$ for the expression of two of the coordinates in terms of the parameters.

In the case of a curve, we have between the coordinates $(x,y,z)$ a twofold relation: two equations $f(x,y,z)=0$, $\phi(x,y,z)=0$ give such a relation; i.e., the curve is here considered as the intersection of two surfaces (but the curve is not always the complete intersection of two surfaces, and there are hence difficulties); or, again, the coordinates may be given each of them as a function of a single variable parameter. The form $y=\phi x$, $z=\psi x$, where two of the coordinates are given in terms of the third, is a particular case of each of these modes of representation.

\textbf{29.} The remarks under plane geometry as to descriptive and metrical propositions, and as to the non-metrical character of the method of coordinates when used for the proof of a descriptive proposition, apply also to solid geometry; and they might be illustrated in like manner by the instance of the theorem of the radical centre of four spheres. The proof is obtained from the consideration that $S$ and $S'$ being each of them a function of the form $x^2+y^2+z^2+ax+by+cz+d$, the difference $S-S'$ is a mere linear function of the coordinates, and consequently that $S-S'=0$ is the equation of the plane containing the circle of intersection of the two spheres $S=0$ and $S'=0$.

\newpage
\noindent\textsc{Page~570.} \subsection*{Metrical Theory}
\addcontentsline{toc}{subsection}{Metrical Theory (Solid Geometry)}

\textbf{30.} The foundation in solid geometry of the metrical theory is, in fact, the before-mentioned theorem that, if a finite right line $PQ$ be projected upon any other line $OO'$ by lines perpendicular to $OO'$, then the length of the projection $P'Q'$ is equal to the length of $PQ$ multiplied by the cosine of its inclination to $P'Q'$; or (in the form in which it is now convenient to state the theorem) the perpendicular distance $PQ'$ of two parallel planes is equal to the inclined distance $PQ$ multiplied by the cosine of the inclination. Hence also the algebraical sum of the projections of the sides of a closed polygon upon any line is $=0$; or, reversing the signs of certain sides and considering the polygon as made up of two broken lines each extending from the same initial to the same terminal point, the sum of the projections of the one set of lines upon any line is equal to the sum of the projections of the other set of lines upon the same line. When any of the lines are at right angles to the given line (or, what is the same thing, in a plane at right angles to the given line), the projections of these lines severally vanish.

\textbf{31.} Consider the skew quadrilateral $QMNP$, the sides $QM$, $MN$, $NP$ being respectively parallel to the three rectangular axes $Ox$, $Oy$, $Oz$; let the lengths of these sides be $\xi$, $\eta$, $\zeta$, and that of the side $QP$ be $=p$; and let the cosines of the inclinations (or say the cosine-inclinations) of $p$ to the three axes be $\alpha$, $\beta$, $\gamma$; then projecting successively on the three sides and on $QP$, we have
\[
\xi=p\alpha,\quad \eta=p\beta,\quad \zeta=p\gamma,
\]
and
\[
p=\xi\alpha+\eta\beta+\zeta\gamma,
\]
whence $p^2=\xi^2+\eta^2+\zeta^2$, which is the relation between a distance $p$ and its projections $\xi$, $\eta$, $\zeta$ upon three rectangular axes. And from the same equations we obtain $\alpha^2+\beta^2+\gamma^2=1$, which is a relation connecting the cosine-inclinations of a line to three rectangular axes.

Suppose we have through $Q$ any other line $QT$, and let the cosine-inclinations of this to the axes be $\alpha'$, $\beta'$, $\gamma'$, and $\theta$ be its cosine-inclination to $QP$; also let $\varpi$ be the length of the projection of $QP$ upon $QT$; then projecting on $QT$, we have
\[
\varpi = \alpha'\xi + \beta'\eta + \gamma'\zeta = p\theta.
\]
And in the last equation substituting for $\xi$, $\eta$, $\zeta$ their values $p\alpha$, $p\beta$, $p\gamma$, we find
\[
\theta = \alpha\alpha' + \beta\beta' + \gamma\gamma',
\]
which is an expression for the mutual cosine-inclination of two lines, the cosine-inclinations of which to the axes are $\alpha$, $\beta$, $\gamma$ and $\alpha'$, $\beta'$, $\gamma'$ respectively. We have of course $\alpha^2+\beta^2+\gamma^2=1$, and $\alpha'^2+\beta'^2+\gamma'^2=1$, and hence also
\[
1-\theta^2 = (\beta\gamma'-\beta'\gamma)^2 + (\gamma\alpha'-\gamma'\alpha)^2 + (\alpha\beta'-\alpha'\beta)^2,
\]
so that the sine of the inclination can only be expressed as a square root. These formulae are the foundation of spherical trigonometry.

\newpage
\noindent\textsc{Page~571.} \subsection*{The Line, Plane, and Sphere}
\addcontentsline{toc}{subsection}{The Line, Plane, and Sphere}

\textbf{32.} The foregoing formulae give at once the equations of these loci.

For first, taking $Q$ to be a fixed point, coordinates $(a,b,c)$ and the cosine-inclinations $(\alpha,\beta,\gamma)$ to be constant, then $P$ will be a point in the line through $Q$ in the direction thus determined; or, taking $(x,y,z)$ for its coordinates, these will be the current coordinates of a point in the line. The values of $\xi$, $\eta$, $\zeta$ then are $x-a$, $y-b$, $z-c$, and we thus have
\[
\frac{x-a}{\alpha} = \frac{y-b}{\beta} = \frac{z-c}{\gamma} (= p),
\]
which (omitting the last equation, $=p$) are the equations of the line through the given point $(a,b,c)$, the cosine-inclinations to the axes being $\alpha$, $\beta$, $\gamma$, where $\alpha^2+\beta^2+\gamma^2=1$. The equations may be written in a form not involving the cosine-inclinations; viz. they may be written
\[
\frac{x-a}{l} = \frac{y-b}{m} = \frac{z-c}{n},
\]
and then $l$, $m$, $n$, instead of being equal, will only be proportional to the cosine-inclinations. Using the last equation, and writing
\[
x, y, z = a+\alpha p,\ b+\beta p,\ c+\gamma p,
\]
these are expressions for the current coordinates in terms of a parameter $p$, which is in fact the distance from the fixed point $(a,b,c)$.

It is easy to see that, if the coordinates $(x,y,z)$ are connected by any two linear equations, these equations can always be brought into the foregoing form, and hence that the two linear equations represent a line.

Secondly, taking for greater simplicity the point $Q$ to be coincident with the origin, and $\alpha'$, $\beta'$, $\gamma'$, $p$ to be constant, then $p$ is the perpendicular distance of a plane from the origin, and $\alpha'$, $\beta'$, $\gamma'$ are the cosine-inclinations of this distance to the axes ($\alpha'^2+\beta'^2+\gamma'^2=1$). Now $P$ is any point in this plane; taking its coordinates to be $(x,y,z)$, then $(\xi,\eta,\zeta)$ are $=(x,y,z)$, and the foregoing equation $p=\alpha'\xi+\beta'\eta+\gamma'\zeta$ becomes
\[
\alpha' x+\beta' y+\gamma' z=p,
\]
which is the equation of the plane in question.

If, more generally, $Q$ is not coincident with the origin, then, taking its coordinates to be $(a,b,c)$, and writing $p_1$ instead of $p$, the equation is
\[
\alpha'(x-a)+\beta'(y-b)+\gamma'(z-c)=p_1;
\]
and we thence have $p_1=p-(a\alpha'+b\beta'+c\gamma')$, which is an expression for the perpendicular distance of the point $(a,b,c)$ from the plane in question.

It is obvious that any linear equation $Ax+By+Cz+D=0$ between the coordinates can always be brought into the foregoing form, and hence that such equation represents a plane.

\newpage
\noindent\textsc{Page~572.} Thirdly, supposing $Q$ to be a fixed point, coordinates $(a,b,c)$ and the distance $QP$, $=p$, to be constant, say this is $=d$, then, as before, the values of $\xi$, $\eta$, $\zeta$ are $x-a$, $y-b$, $z-c$, and the equation $\xi^2+\eta^2+\zeta^2=p^2$ becomes
\[
(x-a)^2+(y-b)^2+(z-c)^2=d^2,
\]
which is the equation of the sphere, coordinates of the centre $=(a,b,c)$ and radius $=d$.

A quadric equation wherein the terms of the second order are $x^2+y^2+z^2$, viz.\ an equation
\[
x^2+y^2+z^2+Ax+By+Cz+D=0,
\]
can always, it is clear, be brought into the foregoing form; and it thus appears that this is the equation of a sphere, coordinates of the centre $=-\tfrac12A$, $-\tfrac12B$, $-\tfrac12C$, and squared radius $=\tfrac14(A^2+B^2+C^2)-D$.

\subsection*{Cylinders, Cones, Ruled Surfaces}
\addcontentsline{toc}{subsection}{Cylinders, Cones, Ruled Surfaces}

\textbf{33.} A singly infinite system of lines, or a system of lines depending upon one variable parameter, forms a surface; and the equation of the surface is obtained by eliminating the parameter between the two equations of the line.

If the lines all pass through a given point, then the surface is a cone; and, in particular, if the lines are all parallel to a given line, then the surface is a cylinder.

Beginning with this last case, suppose the lines are parallel to the line $x=mz$, $y=nz$, the equations of a line of the system are $x=mz+a$, $y=nz+b$, where $a$, $b$ are supposed to be functions of the variable parameter, or, what is the same thing, there is between them a relation $f(a,b)=0$: we have $a=x-mz$, $b=y-nz$, and the result of the elimination of the parameter therefore is $f(x-mz, y-nz)=0$, which is thus the general equation of the cylinder the generating lines whereof are parallel to the line $x=mz$, $y=nz$. The equation of the section by the plane $z=0$ is $f(x,y)=0$, and conversely if the cylinder be determined by means of its curve of intersection with the plane $z=0$, then, taking the equation of this curve to be $f(x,y)=0$, the equation of the cylinder is $f(x-mz, y-nz)=0$. Thus, if the curve of intersection be the circle $(x-\alpha)^2+(y-\beta)^2=\gamma^2$, we have $(x-mz-\alpha)^2+(y-nz-\beta)^2=\gamma^2$ as the equation of an oblique cylinder on this base, and thus also $(x-\alpha)^2+(y-\beta)^2=\gamma^2$ as the equation of the right cylinder.

If the lines all pass through a given point $(a,b,c)$, then the equations of a line are $x-a=\alpha(z-c)$, $y-b=\beta(z-c)$, where $\alpha$, $\beta$ are functions of the variable parameter, or, what is the same thing, there exists between them an equation $f(\alpha,\beta)=0$; the elimination of the parameter gives, therefore, $f\left(\dfrac{x-a}{z-c}, \dfrac{y-b}{z-c}\right)=0$; and this equation, or, what is the same thing, any homogeneous equation $f(x-a, y-b, z-c)=0$, or, taking $f$ to be a rational and integral function of the order $n$, say $(*\,\raisebox{0.5ex}{\text{$\!(x-a,y-b,z-c)$}})^n=0$, is the general equation of the cone having the point $(a,b,c)$ for its vertex. Taking the vertex to be at the origin, the equation is $(*\,\raisebox{0.5ex}{\text{$\!(x,y,z)$}})^n=0$; and, in particular, $(*\,\raisebox{0.5ex}{\text{$\!(x,y,z)$}})^2=0$ is the equation of a cone of the second order, or quadricone, having the origin for its vertex.

\newpage
\noindent\textsc{Page~573.} \textbf{34.} In the general case of a singly infinite system of lines, the locus is a ruled surface (or \textit{regulus}). If the system be such that a line does not intersect the consecutive line, then the surface is a \textit{skew surface}, or \textit{scroll}; but if it be such that each line intersects the consecutive line, then it is a \textit{developable}, or \textit{torse}.

Suppose, for instance, that the equations of a line (depending on the variable parameter $\theta$) are
\[
\frac{x}{a}=\frac{y-\theta b}{\theta b}=\frac{z-c}{-c},\quad\text{or say}\quad
\frac{x}{a}=\frac{y-\theta b}{\theta b}=\frac{z-c}{-c},
\]
then, eliminating $\theta$, we have $\dfrac{x^2}{a^2}+\dfrac{y^2}{b^2}-\dfrac{z^2}{c^2}=1$, if $z\neq c$, or say $x^2/a^2+y^2/b^2-z^2/c^2=1$, the equation of a quadric surface, afterwards called the hyperboloid of one sheet; this surface is consequently a scroll. It is to be remarked that we have upon the surface a second singly infinite series of lines; the equations of a line of this second system (depending on the variable parameter $\phi$) are
\[
\frac{x}{a}-\frac{z}{c}=\lambda\left(1-\frac{y}{b}\right),\quad
\frac{x}{a}+\frac{z}{c}=\frac{1}{\lambda}\left(1+\frac{y}{b}\right).
\]
It is easily shown that any line of the one system intersects every line of the other system.

Considering any curve (of double curvature) whatever, the tangent lines of the curve form a singly infinite system of lines, each line intersecting the consecutive line of the system, that is, they form a developable, or torse; the curve and torse are thus inseparably connected together, forming a single geometrical figure. A plane through three consecutive points of the curve (or osculating plane of the curve) contains two consecutive tangents, that is, two consecutive lines of the torse, and is thus a tangent plane of the torse along a generating line.

\subsection*{Transformation of Coordinates}
\addcontentsline{toc}{subsection}{Transformation of Coordinates (Solid Geometry)}

\textbf{35.} There is no difficulty in changing the origin, and it is for brevity assumed that the origin remains unaltered. We have, then, two sets of rectangular axes, $Ox$, $Oy$, $Oz$, and $Ox_1$, $Oy_1$, $Oz_1$, the mutual cosine-inclinations being shown by the diagram
\begin{center}
\begin{tabular}{c|ccc}
& $x$ & $y$ & $z$ \\
\hline
$x_1$ & $\alpha$ & $\beta$ & $\gamma$ \\
$y_1$ & $\alpha'$ & $\beta'$ & $\gamma'$ \\
$z_1$ & $\alpha''$ & $\beta''$ & $\gamma''$ \\
\end{tabular}
\end{center}
that is, $\alpha$, $\beta$, $\gamma$ are the cosine-inclinations of $Ox_1$ to $Ox$, $Oy$, $Oz$; $\alpha'$, $\beta'$, $\gamma'$ those of $Oy_1$, \&c.

\newpage
\noindent\textsc{Page~574.} And this diagram gives also the linear expressions of the coordinates $(x_1,y_1,z_1)$ or $(x,y,z)$ of either set in terms of those of the other set; we thus have
\begin{align*}
x_1 &= \alpha x + \beta y + \gamma z, & x &= \alpha x_1 + \alpha' y_1 + \alpha'' z_1,\\
y_1 &= \alpha' x + \beta' y + \gamma' z, & y &= \beta x_1 + \beta' y_1 + \beta'' z_1,\\
z_1 &= \alpha'' x + \beta'' y + \gamma'' z, & z &= \gamma x_1 + \gamma' y_1 + \gamma'' z_1,
\end{align*}
which are obtained by projection, as above explained. Each of these equations is, in fact, nothing else than the before-mentioned equation $p=\alpha'\xi+\beta'\eta+\gamma'\zeta$, adapted to the problem in hand.

But we have to consider the relations between the nine coefficients. By what precedes, or by the consideration that we must have identically $x^2+y^2+z^2=x_1^2+y_1^2+z_1^2$, it appears that these satisfy the relations
\begin{gather*}
\alpha^2+\beta^2+\gamma^2=1,\quad \alpha'^2+\beta'^2+\gamma'^2=1,\quad \alpha''^2+\beta''^2+\gamma''^2=1,\\
\alpha'\alpha''+\beta'\beta''+\gamma'\gamma''=0,\quad
\alpha''\alpha+\beta''\beta+\gamma''\gamma=0,\quad
\alpha\alpha'+\beta\beta'+\gamma\gamma'=0,
\end{gather*}
either set of six equations being implied in the other set.

It follows that the square of the determinant
\[
\begin{vmatrix}
\alpha & \beta & \gamma \\
\alpha' & \beta' & \gamma' \\
\alpha'' & \beta'' & \gamma''
\end{vmatrix}
\]
is $=1$; and hence that the determinant itself is $=\pm 1$. The distinction of the two cases is an important one: if the determinant is $=+1$, then the axes $Ox_1$, $Oy_1$, $Oz_1$ are such that they can by a rotation about $O$ be brought to coincide with $Ox$, $Oy$, $Oz$ respectively; if it is $=-1$, then they cannot. But in the latter case, by measuring $x_1$, $y_1$, $z_1$ in the opposite directions we change the signs of all the coefficients and so make the determinant to be $=+1$; hence this case need alone be considered, and it is accordingly assumed that the determinant is $=+1$. This being so, it is found that we have a further set of nine equations, $\alpha=\beta'\gamma''-\beta''\gamma'$, \&c.; that is, the coefficients arranged as in the diagram have the values
\begin{center}
\begin{tabular}{c|ccc}
& $x$ & $y$ & $z$ \\
\hline
$x_1$ & $\beta'\gamma''-\beta''\gamma'$ & $\gamma'\alpha''-\gamma''\alpha'$ & $\alpha'\beta''-\alpha''\beta'$ \\
$y_1$ & $\beta''\gamma-\beta\gamma''$ & $\gamma''\alpha-\gamma\alpha''$ & $\alpha''\beta-\alpha\beta''$ \\
$z_1$ & $\beta\gamma'-\beta'\gamma$ & $\gamma\alpha'-\gamma'\alpha$ & $\alpha\beta'-\alpha'\beta$ \\
\end{tabular}
\end{center}

\newpage
\noindent\textsc{Page~575.} \textbf{36.} It is important to express the nine coefficients in terms of three independent quantities. A solution which, although unsymmetrical, is very convenient in Astronomy and Dynamics is to use for the purpose the three angles $\theta$, $\phi$, $\tau$ of fig.~19; say $\theta=$ longitude of the node; $\phi=$ inclination; and $\tau=$ longitude of $x_1$ (from node, $\Upsilon$).

The diagram of transformation then is
\begin{center}
\begin{tabular}{c|ccc}
& $x$ & $y$ & $z$ \\
\hline
$x_1$ & $\cos\tau\cos\theta-\sin\tau\sin\theta\cos\phi$ & $\cos\tau\sin\theta+\sin\tau\cos\theta\cos\phi$ & $\sin\tau\sin\phi$ \\
$y_1$ & $-\sin\tau\cos\theta-\cos\tau\sin\theta\cos\phi$ & $-\sin\tau\sin\theta+\cos\tau\cos\theta\cos\phi$ & $\cos\tau\sin\phi$ \\
$z_1$ & $\sin\theta\sin\phi$ & $-\cos\theta\sin\phi$ & $\cos\phi$ \\
\end{tabular}
\end{center}

But a more elegant solution (due to Rodrigues) is that contained in the diagram
\begin{center}
\begin{tabular}{c|ccc}
& $x$ & $y$ & $z$ \\
\hline
$x_1$ & $1+\lambda^2-\mu^2-\nu^2$ & $2(\lambda\mu-\nu)$ & $2(\lambda\nu+\mu)$ \\
$y_1$ & $2(\lambda\mu+\nu)$ & $1-\lambda^2+\mu^2-\nu^2$ & $2(\mu\nu-\lambda)$ \\
$z_1$ & $2(\nu\lambda-\mu)$ & $2(\mu\nu+\lambda)$ & $1-\lambda^2-\mu^2+\nu^2$ \\
\end{tabular}
\end{center}
$\div(1+\lambda^2+\mu^2+\nu^2)$.

The nine coefficients of transformation are the nine functions of the diagram, each divided by $1+\lambda^2+\mu^2+\nu^2$; the expressions contain as they should do the three arbitrary quantities $\lambda$, $\mu$, $\nu$; and the identity $x_1^2+y_1^2+z_1^2=x^2+y^2+z^2$ can be at once verified.

It may be added that the transformation can be expressed in the quaternion form
\[
ix_1+jy_1+kz_1=(1+\Lambda)(ix+jy+kz)(1+\Lambda)^{-1},
\]
where $\Lambda$ denotes the vector $i\lambda+j\mu+k\nu$.

\newpage
\noindent\textsc{Page~576.} \subsection*{Quadric Surfaces (Paraboloids, Ellipsoid, Hyperboloids)}
\addcontentsline{toc}{subsection}{Quadric Surfaces}

\textbf{37.} It appears, by a discussion of the general equation of the second order $(a,\ldots\,\raisebox{0.5ex}{\text{$\!(x,y,z,1)$}})^2=0$, that the proper quadric surfaces\footnote{The improper quadric surfaces represented by the general equation of the second order are (1) the pair of planes or plane-pair, including as a special case the twice repeated plane, and (2) the cone, including as a special case the cylinder. There is but one form of cone; but the cylinder may be parabolic, elliptic, or hyperbolic.} represented by such an equation are the following five surfaces ($a$ and $b$ positive):
\begin{align*}
\text{(1)}\quad &\frac{x^2}{a^2}+\frac{y^2}{b^2}=2z, && \text{elliptic paraboloid,}\\
\text{(2)}\quad &\frac{x^2}{a^2}-\frac{y^2}{b^2}=2z, && \text{hyperbolic paraboloid,}\\
\text{(3)}\quad &\frac{x^2}{a^2}+\frac{y^2}{b^2}+\frac{z^2}{c^2}=1, && \text{ellipsoid,}\\
\text{(4)}\quad &\frac{x^2}{a^2}+\frac{y^2}{b^2}-\frac{z^2}{c^2}=1, && \text{hyperboloid of one sheet,}\\
\text{(5)}\quad &\frac{x^2}{a^2}+\frac{y^2}{b^2}-\frac{z^2}{c^2}=-1, && \text{hyperboloid of two sheets.}
\end{align*}
It is at once seen that these are distinct surfaces; and the equations also show very readily the general form and mode of generation of the several surfaces.

In the elliptic paraboloid (fig.~20), the sections by the planes of $zx$ and $zy$ are the parabolas
\[
\frac{x^2}{2a^2}=z,\quad \frac{y^2}{2b^2}=z
\]
having the common axis $Oz$; and the section by any plane $z=\gamma$ parallel to that of $xy$ is the ellipse
\[
\frac{x^2}{a^2}+\frac{y^2}{b^2}=2\gamma,
\]
so that the surface is generated by a variable ellipse moving parallel to itself along the parabolas as directrices.

\newpage
\noindent\textsc{Page~577.} In the hyperbolic paraboloid (fig.~21), the sections by the planes of $zx$, $zy$ are the parabolas
\[
\frac{x^2}{2a^2}=z,\quad -\frac{y^2}{2b^2}=z
\]
having the opposite axes $Oz$, $Oz'$; and the section by a plane $z=\gamma$ parallel to that of $xy$ is the hyperbola
\[
\frac{x^2}{a^2}-\frac{y^2}{b^2}=2\gamma,
\]
which has its transverse axis parallel to $Ox$ or $Oy$ according as $\gamma$ is positive or negative. The surface is thus generated by a variable hyperbola moving parallel to itself along the parabolas as directrices. The form is best seen from fig.~22, which represents the sections by planes parallel to the plane of $xy$, or say the contour lines; the continuous lines are the sections above the plane of $xy$, and the dotted lines the sections below this plane. The form is, in fact, that of a saddle.

In the ellipsoid (fig.~23), the sections by the planes of $zx$, $zy$, and $xy$ are each of them an ellipse, and the section by any parallel plane is also an ellipse. The surface may be considered as generated by an ellipse moving parallel to itself along two ellipses as directrices.

In the hyperboloid of one sheet (fig.~24), the sections by the planes of $zx$, $zy$ are the hyperbolas
\[
\frac{x^2}{a^2}-\frac{z^2}{c^2}=1,\quad \frac{y^2}{b^2}-\frac{z^2}{c^2}=1
\]
having a common conjugate axis $zOz'$; the section by the plane of $xy$, and that by any parallel plane, is an ellipse; and the surface may be considered as generated by a variable ellipse moving parallel to itself along the two hyperbolas as directrices.

In the hyperboloid of two sheets (fig.~25), the sections by the planes of $zx$ and $zy$ are the hyperbolas
\[
\frac{x^2}{a^2}-\frac{z^2}{c^2}=1,\quad \frac{y^2}{b^2}-\frac{z^2}{c^2}=1
\]
having the common transverse axis $zOz'$; the section by any plane $z=\pm\gamma$ parallel to that of $xy$, $\gamma$ being in absolute magnitude $>c$, is the ellipse
\[
\frac{x^2}{a^2}+\frac{y^2}{b^2}=\frac{\gamma^2}{c^2}-1,
\]
and the surface, consisting of two distinct portions or sheets, may be considered as generated by a variable ellipse moving parallel to itself along the hyperbolas as directrices.

The hyperbolic paraboloid is such (and it is easy from the figure to understand how this may be the case) that there exist upon it two singly infinite series of right lines. The same is the case with the hyperboloid of one sheet (ruled or skew hyperboloid, as with reference to this property it is termed). If we imagine two equal and parallel circular disks, their points connected by strings of equal length, so that these are the generating lines of a right circular cylinder, then by turning one of the disks about its centre through the same angle in one or the other direction, the strings will in each case generate one and the same hyperboloid, and will in regard to it be the two systems of lines on the surface, or say the two systems of generating lines; and the general configuration is the same when instead of circles we have ellipses. It has been already shown analytically that the equation $\dfrac{x^2}{a^2}+\dfrac{y^2}{b^2}-\dfrac{z^2}{c^2}=1$ is satisfied by each of two pairs of linear relations between the coordinates.

\newpage
\noindent\textsc{Page~579.} \subsection*{Curves; Tangent, Osculating Plane, Curvature, \&c.}
\addcontentsline{toc}{subsection}{Curves; Tangent, Osculating Plane, Curvature, \&c.}

\textbf{38.} It will be convenient to consider the coordinates $(x,y,z)$ of the point on the curve as given in terms of a parameter $\theta$, so that $dx$, $dy$, $dz$, $d^2x$, \&c., will be proportional to $\dfrac{dx}{d\theta}$, $\dfrac{dy}{d\theta}$, $\dfrac{dz}{d\theta}$, $\dfrac{d^2x}{d\theta^2}$, \&c.; and in them $\xi$, $\eta$, $\zeta$ are used as current coordinates.

The tangent is the line through the point $(x,y,z)$ and the consecutive point $(x+dx, y+dy, z+dz)$; its equations therefore are
\[
\frac{\xi-x}{dx}=\frac{\eta-y}{dy}=\frac{\zeta-z}{dz}.
\]
The osculating plane is the plane through the point and two consecutive points, and contains therefore the tangent; its equation is
\[
\begin{vmatrix}
\xi-x & \eta-y & \zeta-z \\
dx & dy & dz \\
d^2x & d^2y & d^2z
\end{vmatrix}=0,
\]
or, what is the same thing,
\begin{multline}
(\xi-x)(dy\,d^2z-dz\,d^2y)+(\eta-y)(dz\,d^2x-dx\,d^2z)\\
+(\zeta-z)(dx\,d^2y-dy\,d^2x)=0.
\end{multline}
The normal plane is the plane through the point at right angles to the tangent. It meets the osculating plane in a line called the \textit{principal normal}; and drawing through the point a line at right angles to the osculating plane, this is called the \textit{binormal}. We have thus at the point a set of three rectangular axes---the tangent, the principal normal, and the binormal.

We have through the point and three consecutive points a sphere of spherical curvature, the centre and radius thereof being the centre, and radius, of spherical curvature. The sphere is met by the osculating plane in the circle of absolute curvature, the centre and radius thereof being the centre, and radius, of absolute curvature. The centre of absolute curvature is also the intersection of the principal normal by the normal plane at the consecutive point.

\newpage
\noindent\textsc{Page~580.} \subsection*{Surfaces; Tangent Lines and Plane, Curvature, \&c.}
\addcontentsline{toc}{subsection}{Surfaces; Tangent Lines and Plane, Curvature, \&c.}

\textbf{39.} It will be convenient to consider the surface as given by an equation $f(x,y,z)=0$ between the coordinates; taking $(x,y,z)$ for the coordinates of a given point, and $(x+dx, y+dy, z+dz)$ for those of a consecutive point, the increments $dx$, $dy$, $dz$ satisfy the condition
\[
\frac{df}{dx}dx+\frac{df}{dy}dy+\frac{df}{dz}dz=0;
\]
but the ratio of two of the increments, suppose $dx:dy$, may be regarded as arbitrary.

Only a part of the analytical formulae will be given; in them $\xi$, $\eta$, $\zeta$ are used as current coordinates.

We have through the point a singly infinite series of right lines, each meeting the surface in a consecutive point, or say having each of them two-point intersection with the surface. These lines lie all of them in a plane which is the tangent plane; its equation is
\[
(\xi-x)\frac{df}{dx}+(\eta-y)\frac{df}{dy}+(\zeta-z)\frac{df}{dz}=0,
\]
as is at once verified by observing that this equation is satisfied (irrespectively of the value of $dx:dy$) on writing therein $\xi$, $\eta$, $\zeta=x+dx$, $y+dy$, $z+dz$.

\newpage
\noindent\textsc{Page~581.} The line through the point at right angles to the tangent plane is called the \textit{normal}; its equations are
\[
\frac{\xi-x}{\dfrac{df}{dx}}=\frac{\eta-y}{\dfrac{df}{dy}}=\frac{\zeta-z}{\dfrac{df}{dz}}.
\]
In the series of tangent lines there are in general two (real or imaginary) lines, each of which meets the surface in a second consecutive point, or say it has three-point intersection with the surface; these are called the \textit{chief-tangents} (Haupt-tangenten). The tangent-plane cuts the surface in a curve, having at the point of contact a node (double point), the tangents to the two branches being the chief-tangents.

In the case of a quadric surface the curve of intersection, qua curve of the second order, can only have a node by breaking up into a pair of lines; that is, every tangent-plane meets the surface in a pair of lines, or we have on the surface two singly infinite systems of lines; these are real for the hyperbolic paraboloid and the hyperboloid of one sheet, imaginary in other cases.

At each point of a surface the chief-tangents determine two directions; and passing along one of them to a consecutive point, and thence (without abrupt change of direction) along the new chief-tangent to a consecutive point, and so on, we have on the surface a chief-tangent curve; and there are, it is clear, two singly infinite series of such curves. In the case of a quadric surface, the curves are the right lines on the surface.

\textbf{40.} If at the point we draw in the tangent-plane two lines bisecting the angles between the chief-tangents, these lines (which are at right angles to each other) are called the \textit{principal tangents}\footnote{The point on the surface may be such that the directions of the principal tangents become arbitrary; the point is then an \textit{umbilicus}. It is in the text assumed that the point on the surface is not an umbilicus.}. We have thus at each point of the surface a set of rectangular axes, the normal and the two principal tangents.

Proceeding from the point along a principal tangent to a consecutive point on the surface, and thence (without abrupt change of direction) along the new principal tangent to a consecutive point, and so on, we have on the surface a \textit{curve of curvature}; there are, it is clear, two singly infinite series of such curves, cutting each other at right angles at each point of the surface.

Passing from the given point in an arbitrary direction to a consecutive point on the surface, the normal at the given point is not intersected by the normal at the consecutive point; but passing to the consecutive point along a curve of curvature (or, what is the same thing, along a principal tangent) the normal at the given point is intersected by the normal at the consecutive point; we have thus on the normal two centres of curvature, and the distances of these from the point on the surface are the two \textit{principal radii of curvature} of the surface at that point; these are also the radii of curvature of the sections of the surface by planes through the normal and the two principal tangents respectively; or say they are the radii of curvature of the normal sections through the two principal tangents respectively. Take at the point the axis of $z$ in the direction of the normal, and those of $x$ and $y$ in the directions of the principal tangents respectively, then, if the radii of curvature be $a$, $b$ (the signs being such that the coordinates of the two centres of curvature are $z=a$ and $z=b$ respectively), the surface has in the neighbourhood of the point the form of the paraboloid
\[
z=\frac{x^2}{2a}+\frac{y^2}{2b},
\]
and the chief-tangents are determined by the equation $0=\dfrac{x^2}{2a}+\dfrac{y^2}{2b}$. The two centres of curvature may be on the same side of the point or on opposite sides; in the former case $a$ and $b$ have the same sign, the paraboloid is elliptic, and the chief-tangents are imaginary; in the latter case $a$ and $b$ have opposite signs, the paraboloid is hyperbolic, and the chief-tangents are real.

The normal sections of the surface and the paraboloid by the same plane have the same radius of curvature; and it thence readily follows that the radius of curvature of a normal section of the surface by a plane inclined at an angle $\theta$ to that of $zx$ is given by the equation
\[
\frac{1}{\rho}=\frac{\cos^2\theta}{a}+\frac{\sin^2\theta}{b}.
\]
The section in question is that by a plane through the normal and a line in the tangent plane inclined at an angle $\theta$ to the principal tangent along the axis of $x$. To complete the theory, consider the section by a plane having the same trace upon the tangent plane, but inclined to the normal at an angle $\phi$; then it is shown without difficulty (Meunier's theorem) that the radius of curvature of this inclined section of the surface is $=\rho\cos\phi$.

\newpage
\noindent\textsc{Page~583.}

% ============================================================
% PAPER 791: LANDEN
% ============================================================
\section*{791. LANDEN.}
\addcontentsline{toc}{section}{791. Landen}

\begin{flushright}
\small[From the \textit{Encyclop{\ae}dia Britannica}, Ninth Edition, vol.~xiv.~(1882), p.~271.]
\end{flushright}

\textsc{Landen, John}, a distinguished mathematician of the 18th century, was born at Peakirk near Peterborough in Northamptonshire in 1719, and died 15th January 1790 at Milton in the same county. Most of his time was spent in the pursuits of active life, but he early showed a strong talent for mathematical study, which he eagerly cultivated in his leisure hours. In 1762 he was appointed agent to the Earl Fitzwilliam, and held that office to within two years of his death. He lived a very retired life, and saw little or nothing of society; when he did mingle in it, his dogmatism and pugnacity caused him to be generally shunned. He was first known as a mathematician by his essays in the \textit{Ladies' Diary} for 1744. In 1766 he was elected a Fellow of the Royal Society. He was well acquainted and \textit{au courant} with the works of the mathematicians of his own time, and has been called the English D'Alembert. In his Discourse on the ``Residual Analysis,'' in which he proposes to substitute for the method of fluxions a purely algebraical method, he says, ``It is by means of the following theorem, viz.
\[
\frac{X^n-Y^n}{X^m-Y^m}=\frac{X^{n-1}+X^{n-2}Y+\cdots+Y^{n-1}}{X^{m-1}+X^{m-2}Y+\cdots+Y^{m-1}}\quad\text{($n$ and $m$ terms)},
\]
(where $m$ and $n$ are integers), that we are able to perform all the principal operations in our said analysis; and I am not a little surprised that a theorem so obvious, and of such vast use, should so long escape the notice of algebraists.'' The idea is of course a perfectly legitimate one, and may be compared with that of Lagrange's \textit{Calcul des Fonctions}. His memoir (1775) on the rotatory motion of a body contains (as the author was aware) conclusions at variance with those arrived at by D'Alembert and Euler in their researches on the same subject. He reproduces and further develops and defends his own views in his \textit{Mathematical Memoirs}, and in his paper in the \textit{Philosophical Transactions} for 1785. But Landen's capital discovery is that of the theorem known by his name (obtained in its complete form in the memoir of 1775, and reproduced in the first volume of the \textit{Mathematical Memoirs}) for the expression of the arc of an hyperbola in terms of two elliptic arcs. To find this, he integrates a differential equation derived from the equation
\[
\frac{\dot{x}}{\dot{y}}=\frac{y\sqrt{1+g\,t^2}}{t\sqrt{1+g'\,t^2}},
\]
interpreting geometrically in an ingenious and elegant manner three integrals which present themselves. If in the foregoing equation we write $m=1$, $g=k^2$, and instead of $t$ consider the new variable $y=t\sqrt{1-k^2}$, then
\[
\frac{dy}{\sqrt{(1-y^2)(1-k'^2y^2)}}=\frac{dx}{\sqrt{(1-x^2)(1-k^2x^2)}},
\]
which is the form known as Landen's transformation in the theory of elliptic functions: but his investigation does not lead him to obtain the equivalent of the resulting differential equation
\[
\sqrt{1-y^4}\,\frac{dy}{\sqrt{1-y^4}}=\cdots,
\]
due it would appear to Legendre, and which (over and above Landen's own beautiful result) gives importance to the theorem as leading directly to the quadric transformation of an elliptic integral in regard to the modulus.

The list of his writings is as follows: \textit{Ladies' Diary}, various communications, 1744---1760; papers in the \textit{Phil.\ Trans.}, 1754, 1760, 1768, 1771, 1775, 1777, 1785; \textit{Mathematical Lucubrations}, 1755; \textit{Discourse concerning the Residual Analysis}, 1758; \textit{The Residual Analysis}, book~i., 1764; \textit{Animadversions on Dr Stewart's Method of computing the Sun's Distance from the Earth}, 1771; \textit{Mathematical Memoirs}, 1780, 1789.

\newpage
\noindent\textsc{Page~585.}

% ============================================================
% PAPER 792: LOCUS
% ============================================================
\section*{792. LOCUS.}
\addcontentsline{toc}{section}{792. Locus}

\begin{flushright}
\small[From the \textit{Encyclop{\ae}dia Britannica}, Ninth Edition, vol.~xiv.~(1882), pp.~764, 765.]
\end{flushright}

\textsc{Locus}, in Greek $\tau\acute{o}\pi o\varsigma$, a geometrical term, the invention of the notion of which is attributed to Plato. It occurs in such statements as these: the locus of the points which are at the same distance from a fixed point, or of a point which moves so as to be always at the same distance from a fixed point, is a circle; conversely a circle is the locus of the points at the same distance from a fixed point, or of a point moving so as to be always at the same distance from a fixed point; and so, in general, a curve of any given kind is the locus of the points which satisfy, or of a point moving so as always to satisfy, a given condition. The theory of loci is thus identical with that of curves; and it is in fact in this very point of view that a curve is considered in the article Curve, [785]; see that article, and also Geometry (Analytical), [790]. It is only necessary to add that the notion of a locus is useful as regards determinate problems or theorems; thus, to find the centre of the circle circumscribed about a given triangle $ABC$, we see that the circumscribed circle must pass through the two vertices $A$, $B$, and the locus of the centres of the circles which pass through these two points is the straight line at right angles to the side $AB$ at its mid-point; similarly the circumscribed circle must pass through $A$, $C$, and the locus of the centres of the circles through these two points is the line at right angles to the side $AC$ at its mid-point; thus we get the ordinary construction, and also the theorem that the lines at right angles to the sides, at their mid-points respectively, meet in a point. The notion of a locus applies, of course, not only to plane but also to solid geometry. Here the locus of the points satisfying a single (or onefold) condition is a surface; the locus of the points satisfying two conditions (or a twofold condition) is a curve in space, which is in general a twisted curve or curve of double curvature.

\newpage
\noindent\textsc{Page~586.}

% ============================================================
% PAPER 793: MONGE
% ============================================================
\section*{793. MONGE.}
\addcontentsline{toc}{section}{793. Monge}

\begin{flushright}
\small[From the \textit{Encyclop{\ae}dia Britannica}, Ninth Edition, vol.~xvi.~(1883), pp.~738, 739.]
\end{flushright}

\textsc{Monge, Gaspard} (1746--1818), a French mathematician, the inventor of descriptive geometry, was born at Beaune on the 10th May 1746. He was educated first at the college of the Oratorians at Beaune, and then in their college at Lyons, where, at sixteen, the year after he had been learning physics, he was made a teacher of it. Returning to Beaune for a vacation, he made, on a large scale, a plan of the town, inventing the methods of observation and constructing the necessary instruments; the plan was presented to the town, and preserved in their library. An officer of engineers seeing it wrote to recommend Monge to the commandant of the military school at M\'ezi\`eres, and he was received as draftsman and pupil in the practical school attached to that institution; the school itself was of too aristocratic a character to allow of his admission to it. His manual skill was duly appreciated: ``I was a thousand times tempted,'' he said long afterwards, ``to tear up my drawings in disgust at the esteem in which they were held, as if I had been good for nothing better.'' An opportunity, however, presented itself: being required to work out from data supplied to him the ``defilement'' of a proposed fortress (an operation then only performed by a long arithmetical process), Monge, substituting for this a geometrical method, obtained the result so quickly that the commandant at first refused to receive it---the time necessary for the work had not been taken; but upon examination the value of the discovery was recognized, and the method was adopted. And Monge, continuing his researches, arrived at that general method of the application of geometry to the arts of construction which is now called descriptive geometry. But such was the system in France before the Revolution that the officers instructed in the method were strictly forbidden to communicate it even to those engaged in other branches of the public service; and it was not until many years afterwards that an account of it was published. The method consists, as is well known, in the use of the two halves of a sheet of paper to represent say the planes of $xy$ and $xz$ at right angles to each other, and the consequent representation of points, lines, and figures in space by means of their plan and elevation, placed in a determinate relative position.

In 1768 Monge became professor of mathematics, and in 1771 professor of physics, at M\'ezi\`eres; in 1778 he married Madame Horbon, a young widow whom he had previously defended in a very spirited manner from an unfounded charge; in 1780 he was appointed to a chair of hydraulics at the Lyceum in Paris (held by him together with his appointments at M\'ezi\`eres), and was received as a member of the Academy; his intimate friendship with Berthollet began at this time. In 1783, quitting M\'ezi\`eres, he was, on the death of Bezout, appointed examiner of naval candidates. Although pressed by the minister to prepare for them a complete course of mathematics, he declined to do so, on the ground that it would deprive Madame Bezout of her only income, arising from the sale of the works of her late husband; he wrote, however (1786), his \textit{Trait\'e \'el\'ementaire de la Statique}.

Monge contributed (1770--1790) to the \textit{Memoirs of the Academy of Turin}, the \textit{M\'emoires des Savants \'Etrangers} of the Academy of Paris, the \textit{M\'emoires} of the same Academy, and the \textit{Annales de Chimie}, various mathematical and physical papers. Among these may be noticed the memoir ``Sur la th\'eorie des d\'eblais et des remblais'' (M\'em.\ de l'Acad.\ de Paris, 1781), which, while giving a remarkably elegant investigation in regard to the problem of earthwork referred to in the title, establishes in connexion with it his capital discovery of the curves of curvature of a surface. Euler, in his paper on curvature in the Berlin Memoirs for 1760, had considered, not the normals of the surface, but the normals of the plane sections through a particular normal, so that the question of the intersection of successive normals of the surface had never presented itself to him. Monge's memoir just referred to gives the ordinary differential equation of the curves of curvature, and establishes the general theory in a very satisfactory manner; but the application to the interesting particular case of the ellipsoid was first made by him in a later paper in 1795. A memoir in the volume for 1783 relates to the production of water by the combustion of hydrogen; but Monge's results in this matter had been anticipated by Watts and Cavendish.

In 1792, on the creation by the Legislative Assembly of an executive council, Monge accepted the office of minister of the marine, but retained it only until April 1793. When the Committee of Public Safety made an appeal to the savants to assist in producing the mat\'eriel required for the defence of the republic, he applied himself wholly to these operations, and distinguished himself by his indefatigable activity therein; he wrote at this time his \textit{Description de l'art de fabriquer les canons}, and his \textit{Avis aux ouvriers en fer sur la fabrication de l'acier}. He took a very active part in the measures for the establishment of the Normal School (which existed only during the first four months of the year 1795), and of the School for Public Works, afterwards the Polytechnic School, and was at each of them professor of descriptive geometry; his methods in that science were first published in the form in which the shorthand writers took down his lessons given at the Normal School in 1795, and again in 1798--99. In 1796 Monge was sent into Italy with Berthollet and some artists to receive the pictures and statues levied from several Italian towns, and made there the acquaintance of General Bonaparte.

\newpage
\noindent\textsc{Page~588.} Two years afterwards he was sent to Rome on a political mission, which terminated in the establishment, under Massena, of the shortlived Roman republic; and he thence joined the expedition to Egypt, taking part with his friend Berthollet as well in various operations of the war as in the scientific labours of the Egyptian Institute of Sciences and Arts; they accompanied Bonaparte to Syria, and returned with him in 1798 to France. Monge was appointed president of the Egyptian commission, and he resumed his connexion with the Polytechnic School. His later mathematical papers are published (1794--1816) in the \textit{Journal} and the \textit{Correspondance} of the Polytechnic School. On the formation of the Senate he was appointed a member of that body, with an ample provision and the title of count of Pelusium; but on the fall of Napoleon he was deprived of all his honours, and even excluded from the list of members of the reconstituted Institute. He died at Paris on the 28th July 1818.

For further information see B.\ Brisson, \textit{Notice historique sur Gaspard Monge}; Dupin, \textit{Essai historique sur les services et les travaux scientifiques de Gaspard Monge}, Paris, 1819, which contains (pp.~162--166) a list of Monge's memoirs and works; and the biography by Arago (\textit{\OE uvres}, t.~II., 1854).

Monge's various mathematical papers are to a considerable extent reproduced in the \textit{Application de l'Analyse \`a la G\'eom\'etrie}, 4th edition (last revised by the author), Paris, 1819; the pure text of this is reproduced in the 5th edition (revue, corrig\'ee et annot\'ee par M.\ Liouville), Paris, 1850, which contains also Gauss's Memoir, ``Disquisitiones generales circa superficies curvas,'' and some valuable notes by the editor. The other principal separate works are \textit{Trait\'e \'el\'ementaire de la Statique}, $8^{\text{e}}$ Edition, conform\'ee \`a la pr\'ec\'edente, par M.\ Hachette, et suivie d'une Note etc., par M.\ Cauchy, Paris, 1846; and the \textit{G\'eom\'etrie Descriptive} (originating, as mentioned above, in the lessons given at the Normal School). The 4th edition, published shortly after the author's death, seems to have been substantially the same as the 7th (\textit{G\'eom\'etrie Descriptive}, par G.\ Monge, suivie d'une th\'eorie des Ombres et de la Perspective, extraite des papiers de l'auteur, par M.\ Brisson, Paris, 1847).

\newpage
\noindent\textsc{Page~589.}

% ============================================================
% PAPER 794: NUMBERS, PARTITION OF
% ============================================================
\section*{794. NUMBERS, PARTITION OF.}
\addcontentsline{toc}{section}{794. Numbers, Partition of}

\begin{flushright}
\small[From the \textit{Encyclop{\ae}dia Britannica}, Ninth Edition, vol.~xvii.~(1884), p.~614.]
\end{flushright}

This subject, created by Euler, though relating essentially to positive integer numbers, is scarcely regarded as a part of the Theory of Numbers. We consider in it a number as made up by the addition of other numbers: thus the partitions of the successive numbers 1, 2, 3, 4, 5, 6, \&c., are as follows:
\begin{gather*}
1;\\
2,\ 11;\\
3,\ 21,\ 111;\\
4,\ 31,\ 22,\ 211,\ 1111;\\
5,\ 41,\ 32,\ 311,\ 221,\ 2111,\ 11111;\\
6,\ 51,\ 42,\ 411,\ 33,\ 321,\ 3111,\ 222,\ 2211,\ 21111,\ 111111.
\end{gather*}
These are formed each from the preceding ones; thus, to form the partitions of 6 we take first 6; secondly, 5 prefixed to each of the partitions of 1 (that is, 51); thirdly, 4 prefixed to each of the partitions of 2 (that is, 42, 411); fourthly, 3 prefixed to each of the partitions of 3 (that is, 33, 321, 3111); fifthly, 2 prefixed, not to each of the partitions of 4, but only to those partitions which begin with a number not exceeding 2 (that is, 222, 2211, 21111); and lastly, 1 prefixed to all the partitions of 5 which begin with a number not exceeding 1 (that is, 111111); and so in other cases.

The method gives all the partitions of a number, but we may consider different classes of partitions: the partitions into a given number of parts, or into not more than a given number of parts; or the partitions into given parts, either with repetitions or without repetitions, \&c. It is possible, for any particular class of partitions, to obtain methods more or less easy for the formation of the partitions either of a given number or of the successive numbers 1, 2, 3, \&c. And of course in any case, having obtained the partitions, we can count them and so obtain the number ofpartitions.

\newpage
\noindent\textsc{Page~590.} Another method is by Arbogast's rule of the last and the last but one; in fact, taking the value of $a$ to be unity, and, understanding this letter in each term, the rule gives $b$; $c$, $b^2$; $d$, $bc$, $b^3$; $e$, $bd$, $c^2$, $b^2c$, $b^4$, \&c., which, if $b$, $c$, $d$, $e$, \&c., denote 1, 2, 3, 4, \&c., respectively, are the partitions of 1, 2, 3, 4, \&c., respectively.

An important notion is that of \textit{conjugate partitions}. Thus a partition of 6 is 42; writing this in the form
\[
\begin{array}{cc}
** & ** \\
* & \\
* & \\
\end{array}
\]
and summing the columns instead of the lines, we obtain the conjugate partition 2211; evidently, starting from 2211, the conjugate partition is 42. If we form all the partitions of 6 into not more than three parts, these are
\[
6,\ 51,\ 42,\ 33,\ 411,\ 321,\ 222,
\]
and the conjugates are
\[
111111,\ 21111,\ 2211,\ 222,\ 3111,\ 321,\ 33,
\]
where no part is greater than 3; and so, in general, we have the theorem, the number of partitions of $n$ into not more than $k$ parts is equal to the number of partitions of $n$ with no part greater than $k$.

We have for the number of partitions an analytical theory depending on generating functions; thus for the partitions of a number $n$ with the parts 1, 2, 3, 4, 5, \&c., without repetitions, writing down the product
\[
(1+x)(1+x^2)(1+x^3)(1+x^4)(1+x^5)\ldots = 1+x+x^2+2x^3+\ldots +Nx^n+\ldots,
\]
it is clear that, if $x^\alpha$, $x^\beta$, $x^\gamma$,\ldots\ are terms of the series $x$, $x^2$, $x^3$,\ldots\ for which $\alpha+\beta+\gamma+\ldots=n$, then we have in the development of the product a term $x^n$, and hence that, in the term $Nx^n$ of the product, the coefficient $N$ is equal to the number of partitions of $n$ with the parts 1, 2, 3,\ldots, without repetitions; or say that the product is the generating function (G.\ F.) for the number of such partitions. And so in other cases we obtain a generating function.

Thus for the function
\[
\frac{1}{(1-x)(1-x^2)(1-x^3)(1-x^4)\ldots} = 1+x+2x^2+\ldots+Nx^n+\ldots,
\]
observing that any factor $1/(1-x^k)$ is $=1+x^k+x^{2k}+\ldots$, we see that, in the term $Nx^n$, the coefficient $N$ is equal to the number of partitions of $n$, with the parts 1, 2, 3,\ldots, with repetitions.

Introducing another letter $z$, and considering the function
\[
(1+xz)(1+x^2z)(1+x^3z)\ldots = 1+z(x+x^2+\ldots)+\ldots+Nx^nz^k+\ldots
\]
we see that, in the term $Nx^nz^k$ of the development, the coefficient $N$ is equal to the number of partitions of $n$ into $k$ parts, with the parts 1, 2, 3, 4,\ldots, without repetitions. And similarly, considering the function
\[
\frac{1}{(1-xz)(1-x^2z)(1-x^3z)\ldots} = 1+z(x+x^2+\ldots)+\ldots+Nx^nz^k+\ldots
\]
we see that, in the term $Nx^nz^k$ of the development, the coefficient $N$ is equal to the number of partitions of $n$ into $k$ parts, with the parts 1, 2, 3, 4,\ldots, with repetitions.

\newpage
\noindent\textsc{Page~591.} We have such analytical formulae as
\[
\frac{1}{(1-xz)(1-x^2z)(1-x^3z)\ldots} = 1+\frac{zx}{1-x}+\frac{z^2x^2}{(1-x)(1-x^2)}+\ldots,
\]
which lead to theorems in the Partition of Numbers. A remarkable theorem is
\[
(1-x)(1-x^2)(1-x^3)(1-x^4)\ldots = 1-x-x^2+x^5+x^7-x^{12}-x^{15}+\ldots,
\]
where the only terms are those with an exponent $\tfrac12(3n^2\pm n)$ and for each such pair of terms the coefficient is $(-)^n$. The formula shows that, except for numbers of the form $\tfrac12(3n^2\pm n)$, the number of partitions without repetitions into an odd number of parts is equal to the number of partitions without repetitions into an even number of parts, whereas for the excepted numbers these numbers differ by unity. Thus for the number 11, which is not an excepted number, the two sets of partitions are
\begin{align*}
&11,\ 821,\ 731,\ 641,\ 632,\ 542,\\
&10.1,\ 92,\ 83,\ 74,\ 65,\ 5321,
\end{align*}
in each set 6.

We have
\[
(1-x)(1+x)(1+x^2)(1+x^3)(1+x^4)\ldots = 1;
\]
or, as this may be written,
\[
(1+x)(1+x^2)(1+x^3)(1+x^4)\ldots = \frac{1}{1-x} = 1+x+x^2+x^3+\ldots,
\]
showing that a number $n$ can always be made up, and in one way only, with the parts 1, 2, 4, 8,\ldots. The product on the left-hand side may be taken to $k$ terms only: thus if $k=4$, we have
\[
(1+x)(1+x^2)(1+x^4)(1+x^8) = \frac{1-x^{16}}{1-x} = 1+x+x^2+\ldots+x^{15},
\]
that is, any number from 1 to 15 can be made up, and in one way only, with the parts 1, 2, 4, 8; and similarly any number from 1 to $2^k-1$ can be made up, and in one way only, with the parts 1, 2, 4,\ldots, $2^{k-1}$. A like formula is
\[
\frac{(1-x^3)(1-x^9)(1-x^{27})(1-x^{81})}{x(1-x)(x^2-x^3)(x^5-x^9)(x^{14}-x^{27})(x^{41}-x^{81})} = \ldots,
\]
that is,
\[
x^{-40}+1+x\cdot x^{-39}+1+x^2\cdot\ldots\cdot x^{-1}+1+x+x^2+\ldots+x^{39}+x^{40},
\]
showing that any number from $-40$ to $+40$ can be made up, and that in one way only, with the parts 1, 3, 9, 27 taken positively or negatively; and so in general any number from $-\tfrac12(3^k-1)$ to $+\tfrac12(3^k-1)$ can be made up, and that in one way only, with the parts 1, 3, 9,\ldots, $3^{k-1}$ taken positively or negatively.

\newpage
\noindent\textsc{Page~592.}

% ============================================================
% PAPER 795: NUMBERS, THEORY OF
% ============================================================
\section*{795. NUMBERS, THEORY OF.}
\addcontentsline{toc}{section}{795. Numbers, Theory of}

\begin{flushright}
\small[From the \textit{Encyclop{\ae}dia Britannica}, Ninth Edition, vol.~xvii.~(1884), pp.~614--624.]
\end{flushright}

The Theory of Numbers, or higher arithmetic, otherwise arithmology, is a subject which, originating with Euclid, has in modern times, in the hands of Legendre, Gauss, Lejeune-Dirichlet, Kummer, Kronecker, and others, been developed into a most extensive and interesting branch of mathematics. We distinguish between the ordinary (or say the simplex) theory and the various complex theories.

In the ordinary theory we have, in the first instance, positive integer numbers, the unit or unity 1, and the other numbers 2, 3, 4, 5, \&c. We introduce the zero 0, which is a number \textit{sui generis}, and the negative numbers $-1$, $-2$, $-3$, $-4$, \&c., and we have thus the more general notion of integer numbers, 0, $\pm 1$, $\pm 2$, $\pm 3$, \&c.; $+1$ and $-1$ are units or unities. The sum of any two or more numbers is a number; conversely, any number is a sum of two or more parts; but even when the parts are positive a number cannot be, in a determinate manner, represented as a sum of parts. The product of two or more numbers is a number; but (disregarding the unities $+1$, $-1$, which may be introduced as factors at pleasure) it is not conversely true that every number is a product of numbers. A number such as 2, 3, 5, 7, 11, \&c., which is not a product of numbers, is said to be a \textit{prime number}; and a number which is not prime is said to be \textit{composite}. A number other than zero is thus either prime or composite; and we have the theorem that every composite number is, in a determinate way, a product of prime factors.

We have complex theories in which all the foregoing notions (integer, unity, zero, prime, composite) occur; that which first presented itself was the theory with the unit $i$ ($i^2=-1$); we have here complex numbers, $a+bi$, where $a$ and $b$ are in the before-mentioned (ordinary) sense positive or negative integers, not excluding zero; we have the zero $0$, $=0+0i$, and the four units $1$, $-1$, $i$, $-i$. A number other than zero is here either prime or else composite; for instance, 3, 7, 11, are prime numbers, and $5=(2+i)(2-i)$, $9=3\cdot 3$, $13=(3+2i)(3-2i)$, are composite numbers (generally any positive real prime of the form $4n+3$ is prime, but any positive real prime of the form $4n+1$ is a sum of two squares, and is thus composite). And disregarding unit factors we have, as in the ordinary theory, the theorem that every composite number is, in a determinate way, a product of prime factors.

There is, in like manner, a complex theory involving the cube roots of unity---if $\omega$ be an imaginary cube root of unity ($\omega^2+\omega+1=0$), then the integers of this theory are $a+b\omega$ ($a$ and $b$ real positive or negative integers, including zero); a complex theory with the fifth roots of unity---if $\omega$ be an imaginary fifth root of unity ($\omega^4+\omega^3+\omega^2+\omega+1=0$), then the integers of the theory are $a+b\omega+c\omega^2+d\omega^3$ ($a$, $b$, $c$, $d$, real positive or negative integers, including zero); and so on for the roots of the orders 7, 11, 13, 17, 19. In all these theories, or at any rate for the orders 3, 5, 7 (see No.~37, post), we have the foregoing theorem: disregarding unit factors, a number other than zero is either prime or composite, and every composite number is, in a determinate way, a product of prime factors. But coming to the 23rd roots of unity the theorem ceases to be true. Observe that it is a particular case of the theorem that, if $N$ be a prime number, any integer power of $N$ has for factors only the lower powers of $N$,---for instance, $N^3=N\cdot N^2$; there is no other decomposition $N^3=A\cdot B$. This is obviously true in the ordinary theory, and it is true in the complex theories preceding those for the 3rd, 5th, and 7th roots of unity, and probably in those for the other roots preceding the 23rd roots; but it is not true in the theory for the 23rd roots of unity. We have, for instance, 47, a number not decomposable into factors, but $47^2=AB$, is a product of two numbers each of the form $a+b\omega+\ldots+l\omega^{22}$ ($\omega$ a 23rd root). The theorem recovers its validity by the introduction into the theory of Kummer's notion of an ideal number.

The complex theories above referred to would be more accurately described as theories for the complex numbers involving the periods of the roots of unity: the units are the roots either of the equation $x^{p-1}+x^{p-2}+\ldots+x+1=0$ ($p$ a prime number) or of any equation $x^{\frac{p-1}{q}}+\ldots+1=0$ belonging to a factor of the function of the order $p-1$: in particular, this may be the quadric equation for the periods each of $\tfrac12(p-1)$ roots; they are the theories which were first and have been most completely considered, and which led to the notion of an ideal number. But a yet higher generalization which has been made is to consider the complex theory, the units whereof are the roots of any given irreducible equation which has integer numbers for its coefficients.

There is another complex theory the relation of which to the foregoing is not very obvious, viz.\ Galois's theory of the numbers composed with the imaginary roots of an irreducible congruence, $F(x)=0$ (modulus a prime number $p$); the nature of this will be indicated in the sequel.

In any theory, ordinary or complex, we have a first part, which has been termed (but the name seems hardly wide enough) the theory of congruences; a second part, the theory of homogeneous forms: this includes in particular the theory of the binary quadratic forms $(a,b,c)(x,y)^2$; and a third part, comprising those miscellaneous investigations which do not come properly under either of the foregoing heads.

\newpage
\noindent\textsc{Page~594.} \subsection*{Ordinary Theory, First Part}
\addcontentsline{toc}{subsection}{Ordinary Theory, First Part}

\textbf{1.} We are concerned with the integer numbers 0, $\pm 1$, $\pm 2$, $\pm 3$, \&c., or in the first place with the positive integer numbers 1, 2, 3, 4, 5, 6, \&c. Some of these, 1, 2, 3, 5, 7, \&c., are prime, others, $4=2^2$, $6=2\cdot 3$, \&c., are composite; and we have the fundamental theorem that a composite number is expressible, and that in one way only, as a product of prime factors, $N=a^\alpha b^\beta c^\gamma\ldots$ ($a$, $b$, $c$, \ldots\ primes other than 1; $\alpha$, $\beta$, $\gamma$, \ldots\ positive integers).

Gauss makes the proof to depend on the following steps: (i)~the product of two numbers each smaller than a given prime number is not divisible by this number; (ii)~if neither of two numbers is divisible by a given prime number the product is not so divisible; (iii)~the like as regards three or more numbers; (iv)~a composite number cannot be resolved into factors in more than one way.

\textbf{2.} Proofs will in general be only indicated or be altogether omitted, but, as a specimen of the reasoning in regard to whole numbers, the proofs of these fundamental propositions are given at length. (i)~Let $p$ be the prime number, $a$ a number less than $p$, and if possible let there be a number $b$ less than $p$, and such that $ab$ is divisible by $p$; it is further assumed that $b$ is the only number, or, if there is more than one, then that $b$ is the least number having the property in question; $b$ is greater than 1, for $a$ being less than $p$ is not divisible by $p$. Now $p$ as a prime number is not divisible by $b$, but must lie between two consecutive multiples of $b$, say $mb$ and $(m+1)b$; if so, $p-mb$, which is less than $b$, and since neither $p$ nor $mb$ is divisible by $b$, $p-mb$ is not divisible by $b$; but $ab$ and $mb$ are each of them divisible by $b$, therefore their difference $ab-mb=a(p-mb)+m(ab-p)$, since $ab-p$ is divisible by $b$, must also be divisible by $p$, or, writing $p-mb=b'$, we have $ab'$ divisible by $p$, where $b'$ is less than $b$: so that $b$ is not the least number for which $ab$ is divisible by $p$. (ii)~If $a$ and $b$ are neither of them divisible by $p$, then $a$ divided by $p$ leaves a remainder $\alpha$ which is less than $p$, say we have $a=mp+\alpha$; and similarly $b$ divided by $p$ leaves a remainder $\beta$ which is less than $p$, say we have $b=np+\beta$; then
\[
ab=(mp+\alpha)(np+\beta)=(mnp+n\alpha+m\beta)p+\alpha\beta,
\]
and $\alpha\beta$ is not divisible by $p$, therefore $ab$ is not divisible by $p$. (iii)~The like proof applies to the product of three or more factors $a$, $b$, $c$, \ldots\ (iv)~Suppose that the number $N=a^\alpha b^\beta c^\gamma\ldots$ ($a$, $b$, $c$, \ldots\ prime numbers other than 1), is decomposable in some other way into prime factors; we can have no prime factor $p$, other than $a$, $b$, $c$, \ldots, for no such number can divide $a^\alpha b^\beta c^\gamma\ldots$; and we must have each of the numbers $a$, $b$, $c$, \ldots, for if any one of them, suppose $a$, were wanting, the number $N$ would not be divisible by $a$. Hence the new decomposition if it exists must be a decomposition $N=a^{\alpha'}b^{\beta'}c^{\gamma'}\ldots$; and here, if any two corresponding indices, say $\alpha$, $\alpha'$, are different from each other, then one of them, suppose $\alpha$, is the greater, and we have $N\div a^\alpha=b^\beta c^\gamma\ldots=a^{\alpha'-\alpha}b^{\beta'}c^{\gamma'}\ldots$. That is, we have the number $N\div a^\alpha$ expressed in two different ways as a product, the number $a$ being a factor in the one case, but not a factor in the other case. Thus the two exponents cannot be unequal, that is, we must in fact have $\alpha'=\alpha$, and similarly we have $\beta'=\beta$, $\gamma'=\gamma$, \ldots; that is, there is only the one decomposition $N=a^\alpha b^\beta c^\gamma\ldots$.

\newpage
\noindent\textsc{Page~595.} \textbf{3.} The only numbers divisible by a number $N=a^\alpha b^\beta c^\gamma\ldots$ are the numbers $a^{\alpha'}b^{\beta'}c^{\gamma'}\ldots$, where each exponent $\alpha'$ is equal to or greater than the corresponding exponent $\alpha$. And conversely the only numbers which divide $N$ are those of the form $a^{\alpha'}b^{\beta'}c^{\gamma'}\ldots$, where each index $\alpha'$ is at most equal to the corresponding index $\alpha$; and in particular each or any of the indices $\alpha'$ may be $=0$. Again, the least common multiple of two numbers $N=a^\alpha b^\beta c^\gamma\ldots$ and $N'=a^{\alpha'}b^{\beta'}c^{\gamma'}\ldots$ is $a^{\alpha''}b^{\beta''}c^{\gamma''}\ldots$, where each index $\alpha''$ is equal to the largest of the corresponding indices $\alpha$, $\alpha'$;---observe that any one or more of the indices $\alpha$, $\beta$, $\gamma$,\ldots, $\alpha'$, $\beta'$, $\gamma'$,\ldots\ may be $=0$, so that the theorem extends to the case where either of the numbers $N$, $N'$, has prime factors which are not factors of the other number. And so the greatest common measure of two numbers $N=a^\alpha b^\beta c^\gamma\ldots$ and $N'=a^{\alpha'}b^{\beta'}c^{\gamma'}\ldots$ is $a^{\alpha''}b^{\beta''}c^{\gamma''}\ldots$, where each index $\alpha''$ is equal to the least of the corresponding indices $\alpha$ and $\alpha'$.

\textbf{4.} The divisors of $N=a^\alpha b^\beta c^\gamma\ldots$ are the several terms of the product
\[
(1+a+\ldots+a^\alpha)(1+b+\ldots+b^\beta)(1+c+\ldots+c^\gamma)\ldots,
\]
where unity and the number itself are reckoned each of them as a divisor. Hence the number of divisors is $=(\alpha+1)(\beta+1)(\gamma+1)\ldots$, and the sum of the divisors is
\[
\frac{(a^{\alpha+1}-1)(b^{\beta+1}-1)(c^{\gamma+1}-1)\ldots}{(a-1)(b-1)(c-1)\ldots}.
\]

\textbf{5.} In $N=a^\alpha b^\beta c^\gamma\ldots$ the number of integers less than $N$ and prime to it is
\[
\phi(N)=N\left(1-\frac{1}{a}\right)\left(1-\frac{1}{b}\right)\left(1-\frac{1}{c}\right)\ldots
\]
To find the numbers in question write down the series of numbers 1, 2, 3,\ldots, $N$; strike out all the numbers divisible by $a$, then those divisible by $b$, then those divisible by $c$, and so on; there will remain only the numbers prime to $N$. For actually finding the numbers we may of course in striking out those divisible by $b$ disregard the numbers already struck out as divisible by $a$, and in striking out with respect to $c$ disregard the numbers already struck out as divisible by $a$ or $b$, and so on; but in order to count the remaining numbers it is more convenient to ignore the previous strikings out. Suppose, for a moment, there are only two prime factors $a$ and $b$, then the number of terms struck out as divisible by $a$ is $=N\div a$, and the number of terms struck out as divisible by $b$ is $=N\div b$; but then each term divisible by $ab$ will have been twice struck out; the number of these is $=N\div ab$, and thus the number of the remaining terms is $N(1-1/a-1/b+1/ab)=N(1-1/a)(1-1/b)$. By treating in like manner the case of three or more prime factors $a$, $b$, $c$,\ldots\ we arrive at the general theorem. The formula gives $\phi(1)=1$: viz.\ when $N=1$, there is no factor $1$; and it is necessary to consider $\phi(1)$ as being $=1$. The explanation is that $\phi(N)$ properly denotes the number of integers not greater than $N$ and prime to it; so that, when $N=1$, we have 1 an integer not greater than $N$ and prime to it; but in every other case the two definitions agree.

\newpage
\noindent\textsc{Page~596.} \textbf{6.} If $N$, $N'$, are numbers prime to each other, then $\phi(NN')=\phi(N)\phi(N')$, and so also for any number of numbers having no common divisor; in particular,
\[
\phi(a^\alpha b^\beta c^\gamma\ldots)=\phi(a^\alpha)\phi(b^\beta)\phi(c^\gamma)\ldots;\quad
\phi(a^\alpha)=a^\alpha\left(1-\frac{1}{a}\right).
\]
and the theorem is at once verified. We have $N=\sum\phi(N')$, where the summation extends to all the divisors $N'$ of $N$, unity and the number itself being included; thus $15=\phi(15)+\phi(5)+\phi(3)+\phi(1)=8+4+2+1$.

\textbf{7.} The prime factor of the binomial function $x^N-1$ is
\[
[x^N-1]=\frac{(x^N-1)(x^{N/pq}-1)\ldots}{(x^{N/p}-1)(x^{N/q}-1)\ldots},
\]
a rational and integral function of the degree $\phi(N)$; say this is called $[x^N-1]$, and we have $x^N-1=\prod[x^{N'}-1]$, where the product extends to all the divisors $N'$ of $N$, unity and the number included. For instance
\[
[x^6-1]=\frac{(x^6-1)(x-1)}{(x^3-1)(x^2-1)},
\]
and we have
\[
x^6-1=[x^6-1][x^3-1][x^2-1][x-1]=(x^2-x+1)(x^2+x+1)(x+1)(x-1).
\]

\textbf{8. Congruence to a given modulus.} A number $x$ is congruent to 0, to the modulus $N$, $x=0$ (mod.\ $N$), when $x$ is divisible by $N$; two numbers $x$, $y$ are congruent to the modulus $N$, $x=y$ (mod.\ $N$), when their difference $x-y$ divides by $N$, or, what is the same thing, if $x-y=0$ (mod.\ $N$). Observe that, if $xy=0$ (mod.\ $N$), and $x$ be prime to $N$, then $y=0$ (mod.\ $N$).

\textbf{9. Residues to a given modulus.} For a given modulus $N$ we can always find, and that in an infinity of ways, a set of numbers, say residues, such that every number whatever is, to the modulus $N$, congruent to one and only one of these residues. For instance, the residues may be 0, 1, 2, 3,\ldots, $N-1$ (the residue of a given number is here simply the positive remainder of the number when divided by $N$); or, $N$ being odd, the system may be $0$, $\pm 1$, $\pm 2$,\ldots, $\pm\tfrac12(N-1)$, and, $N$ even, $0$, $\pm 1$, $\pm 2$,\ldots, $\pm\tfrac12(N-2)$, $+\tfrac12N$.

\textbf{10. Prime residues to a given modulus.} Considering only the numbers which are prime to a given modulus $N$, we have here a set of $\phi(N)$ numbers, say $\phi(N)$ prime residues, such that every number prime to $N$ is, to the modulus $N$, congruent to one and only one of these prime residues. For instance, the prime residues may be the numbers less than $N$ and prime to it. In particular, if $N$ is a prime number $p$, then the residues may be the $p-1$ numbers, 1, 2, 3,\ldots, $p-1$.

In all that follows, the letter $p$, in the absence of any statement to the contrary, will be used to denote an odd prime other than unity. A theorem for $p$ may hold good for the even prime 2, but it is in general easy to see whether this is so or not.

\newpage
\noindent\textsc{Page~597.} \textbf{11. Fermat's theorem,} $a^{p-1}-1=0$ (mod.\ $p$). The generalized theorem is $a^{\phi(N)}-1=0$ (mod.\ $N$). The proof of the generalized theorem is as easy as that of the original theorem. Consider the series of the $\phi(N)$ numbers $a$, $b$, $c$, \ldots, each less than $N$ and prime to it; let $x$ be any number prime to $N$, then each of the numbers $xa$, $xb$, $xc$,\ldots, is prime to $N$, and no two of them are congruent to the modulus $N$, that is, we cannot have $x(a-b)=0$ (mod.\ $N$); in fact, $x$ is prime to $N$, and the difference $a-b$ of two positive numbers each less than $N$ will be less than $N$. Hence the numbers $xa$, $xb$, $xc$, \ldots, are in a different order congruent to the numbers $a$, $b$, $c$, \ldots; and multiplying together the numbers of each set we have $x^{\phi(N)}abc\ldots=abc\ldots$ (mod.\ $N$), or, since $a$, $b$, $c$, \ldots, are each prime to $N$, and therefore also the product $abc\ldots$ is prime to $N$, we have $x^{\phi(N)}=1$, or say $x^{\phi(N)}-1=0$ (mod.\ $N$).

In particular, if $N$ be a prime number $=p$, then $\phi(N)$ is $=p-1$, and the theorem is $x^{p-1}-1=0$ (mod.\ $p$), $x$ being now any number not divisible by $p$.

\textbf{12.} The general congruence $f(x)=0$ (mod.\ $p$). $f(x)$ is written to denote a rational and integral function with integer coefficients which may without loss of generality be taken to be each of them less than $p$; it is assumed that the coefficient $A$ of the highest power of $x$ is not $=0$. If there is for $x$ an integer value $\alpha$ such that $f(\alpha)=0$ (mod.\ $p$, throughout), then $\alpha$ is said to be a \textit{root} of the congruence $f(x)=0$; we may, it is clear, for $\alpha$ substitute any value whatever $\alpha'=\alpha+kp$, or say any value $\alpha'$ which is $\equiv\alpha$, but such value $\alpha'$ is considered not as a different root but as the same root of the congruence. We have thus $f(\alpha)=0$; and therefore $f(x)=f(x)-f(\alpha)$, $=(x-\alpha)f_1(x)$, where $f_1(x)$ is a function of like form with $f(x)$, that is, with integer coefficients, but of the next inferior order $n-1$. Suppose there is another root $b$ of the congruence, that is, an integer value $b$ such that $f(b)=0$; we have then $(b-\alpha)f_1(b)=0$, and $b-\alpha$ is not $=0$ (for then $b$ would be the same root as $\alpha$). Hence $f_1(b)=0$, and $f(x)=(x-\alpha)[f_1(x)-f_1(b)]=(x-\alpha)(x-b)f_2(x)$, where $f_2(x)$ is an integral function such as $f(x)$, but of the order $n-2$; and so on, that is, if there exist $n$ different (non-congruent) roots of the congruence $f(x)=0$, then $f(x)=A(x-\alpha)(x-b)\ldots(x-k)$, and the congruence may be written $A(x-\alpha)(x-b)\ldots(x-k)=0$. And this cannot be satisfied by any other value $l$; for if so we should have $A(l-\alpha)(l-b)\ldots(l-k)=0$, that is, some one of the congruences $(l-\alpha)=0$, \&c., would have to be satisfied, and $l$ would be the same as one of the roots $\alpha$, $b$, $c$, \ldots, $k$. That is, a congruence of the order $n$ cannot have more than $n$ roots, and if it have precisely $n$ roots $\alpha$, $b$, $c$, \ldots, $k$, then the form is $f(x)=A(x-\alpha)(x-b)\ldots(x-k)=0$.

Observe that a congruence may have equal roots, viz.\ if the form be $f(x)=A(x-\alpha)^\alpha(x-b)^\beta\ldots=0$, then the roots $\alpha$, $b$, \ldots\ are to be counted $\alpha$ times, $\beta$ times, \ldots\ respectively; but clearly the whole number of roots $\alpha+\beta+\ldots$ is at most $=n$.

It is hardly necessary to remark that this theory of a congruence of the order $n$ is precisely analogous to that of an equation of the order $n$, when only real roots are attended to. The theory of the imaginary roots of a congruence will be considered further on (see No.~41).

\newpage
\noindent\textsc{Page~598.} \textbf{13. The linear congruence} $ax=c$ (mod.\ $b$). This is equivalent to the indeterminate equation $ax+by=c$; if $a$ and $b$ are not prime to each other, but have a greatest common measure $q$, this must also divide $c$; supposing the division performed, the equation becomes $a'x+b'y=c'$, where $a'$ and $b'$ are prime to each other, or, what is the same thing, we have the congruence $a'x=c'$ (mod.\ $b'$). This can always be solved, for, if we consider the $b'$ numbers 0, 1, 2,\ldots, $b'-1$, one and only one of these will be $\equiv c'$ (mod.\ $b'$). Multiplying these by any number $a'$ prime to $b'$, and taking the remainders in regard to $b'$, we reproduce in a different order the same series of numbers 0, 1, 2,\ldots, $b'-1$; that is, in the series $a'$, $2a'$, \ldots, $(b'-1)a'$ there will be one and only one term $\equiv c'$ (mod.\ $b'$), or, calling the term in question $\alpha$, we have $x=\alpha$ as the solution of the congruence $a'x=c'$ (mod.\ $b'$); $a'\alpha-c'$ is then a multiple of $b'$, say it is $=-b'\beta$, and the corresponding value of $y$ is $y=\beta$. We may for $\alpha$ write $\alpha+mb'$, $m$ being any positive or negative integer, not excluding zero (but, as already remarked, this is not considered as a distinct solution of the congruence); the corresponding value of $y$ is clearly $\beta-ma'$.

The value of $x$ can be found by a process similar to that for finding the greatest common measure of the two numbers $a'$ and $c'$; this is what is really done in the apparently tentative process which at once presents itself for small numbers, thus $6x=9$ (mod.\ 35), we have $36x=54$, or, rejecting multiples of 35, $x=12$, or, if we please, $x=35m+12$.

In particular, we can always find a number $\xi$ such that $a'\xi=1$ (mod.\ $b'$); and we have then $x=c'\xi$ as the solution of the congruence $a'x=c'$ (mod.\ $b'$). The value of $\xi$ may be written $\xi=\frac{1}{a'}$ (mod.\ $b'$), where $\frac{1}{a'}$ stands for that integer value $\xi$ which satisfies the original congruence $a'\xi=1$ (mod.\ $b'$); and the value of $x$ may then be written $x\equiv c'\cdot\frac{1}{a'}$ (mod.\ $b'$).

Another solution of the linear congruence is given in No.~21.

\textbf{14. Wilson's theorem,} $1\cdot2\cdot3\ldots(p-1)+1=0$ (mod.\ $p$). It has been seen that, for any prime number $p$, the congruence $x^{p-1}-1=0$ (mod.\ $p$) of the order $p-1$ has the $p-1$ roots 1, 2,\ldots, $p-1$; we have therefore
\[
x^{p-1}-1=(x-1)(x-2)\ldots(x-p+1),
\]
or, comparing the terms independent of $x$, it appears that $1\cdot2\cdot3\ldots(p-1)=-1$, that is,
\[
1\cdot2\cdot3\ldots(p-1)+1\equiv0\pmod{p},
\]
---the required theorem. For instance, where $p=5$, then $1\cdot2\cdot3\cdot4+1\equiv0$ (mod.\ 5), and where $p=7$, then $1\cdot2\cdot3\cdot4\cdot5\cdot6+1\equiv0$ (mod.\ 7).

\newpage
\noindent\textsc{Page~599.} \textbf{15.} A proof on wholly different principles may be given. Suppose, to fix the ideas, $p=7$; consider on a circle 7 points, the summits of a regular heptagon, and join these in any manner so as to form a heptagon (7-gon), the number of regular heptagons (convex or stellated) is $=\tfrac12(7-1)$, $=3$, while the number of remaining heptagons must be divisible by 7, for by making any one rotation we have 6 other heptagons successively, each of the 6 sets of 7 being heptagons which can be obtained from it by a rotation of $\tfrac16$ of $360^\circ$, i.e.\ $1\cdot2\cdot3\cdot4\cdot5\cdot6-\tfrac12(7-1)=0$ (mod.\ 7), whence $1\cdot2\cdot3\cdot4\cdot5\cdot6-7+1=0$, or finally $1\cdot2\cdot3\cdot4\cdot5\cdot6+1\equiv0$ (mod.\ 7). It is clear that the proof applies without alteration to the case of any prime number $p$.

If $p$ is not a prime number, then $1\cdot2\cdot3\ldots(p-1)\not\equiv-1$ (mod.\ $p$); hence the theorem shows directly whether a number $p$ is or is not a prime number; but it is not of any practical utility for this purpose.

\textbf{16. Prime roots of a prime number---application to the binomial equation} $x^p-1=0$. Take, for instance, $p=7$. By what precedes, we have
\[
x^7-1=[x^6-1][x^3-1][x^2-1][x-1]=(x^2-x+1)(x^2+x+1)(x+1)(x-1);
\]
and we have
\[
x^7-1=(x-1)(x-2)(x-3)(x-4)(x-5)(x-6)\pmod{7},
\]
whence also
\[
(x^2-x+1)(x^2+x+1)(x+1)(x-1)=(x-1)(x-2)(x-3)(x-4)(x-5)(x-6).
\]
These two decompositions must agree together, and we in fact have
\begin{gather*}
x^2-x+1=(x-3)(x-5),\\
x^2+x+1=(x-2)(x-4),\quad x+1=x-6,\quad x-1=x-1.
\end{gather*}
In particular, we thus have 3, 5, as the roots of the congruence $x^2-x+1=0$, that is, $[x^6-1]=0$, and these roots 3, 5, are not the roots of any other of the congruences $[x^3-1]=0$, $[x^2-1]=0$, $[x-1]=0$; that is, writing $a=3$ or 5 in the series of numbers $a$, $a^2$, $a^3$, $a^4$, $a^5$, $a^6$, we have $a^k$ as the first term which is $\equiv1$ (mod.\ 7); the series in fact are
\begin{gather*}
3,\ 9,\ 27,\ 81,\ 243,\ 729\equiv3,\ 2,\ 6,\ 4,\ 5,\ 1,\\
5,\ 25,\ 125,\ 625,\ 3125,\ 15625\equiv5,\ 4,\ 6,\ 2,\ 3,\ 1.
\end{gather*}
And so, in general, the congruence $x^{p-1}-1=0$ (mod.\ $p$) has the $p-1$ real roots 1, 2, 3,\ldots, $p-1$; hence the congruence $[x^{p-1}-1]=0$, which is of the order $\phi(p-1)$, has this number $\phi(p-1)$ of real roots; and, calling any one of these $g$, then in the series of powers $g$, $g^2$, $g^3$, \ldots, $g^{p-1}$, the first term which is $\equiv1$ (mod.\ $p$) is $g^{p-1}$, that is, we have $g$, $g^2$, $g^3$, \ldots, $g^{p-1}\equiv1$, 2, 3,\ldots, $p-1$ in a different order. Any such number $g$ is said to be a \textit{prime root} of $p$, and the number of prime roots is $\phi(p-1)$, the number of integers less than and prime to $p-1$.

The notion of a prime root was applied by Gauss to the solution of the binomial equation $x^p-1=0$, or, what is the same thing, to the question of the division of the circle (Kreistheilung), see Equation, Nos.~30 and 31, [786]; and, as remarked in the introduction to the present article, the roots or periods of roots of this equation present themselves as the units of a complex theory in the Theory of Numbers.

\newpage
\noindent\textsc{Page~600.} \textbf{17.} Any number $x$ less than $p$ is $=g^m$, and, if $m$ is not prime to $p-1$, but has with it a greatest common measure $e$, suppose $m=ke$, $p-1=e\cdot f$, then
\[
x^f=g^{ke\cdot f}=g^{k(p-1)}\equiv1;
\]
that is, $x^f\equiv1$; and it is easily seen that in the series of powers $x$, $x^2$,\ldots, $x^f$, we have $x^f$ as the first term which is $\equiv1$ (mod.\ $p$). A number $=g^m$, where $m$ is not prime to $p-1$, is thus not a prime root; and it further appears that, $g$ being any particular prime root, the $\phi(p-1)$ prime roots are the numbers $g^m$, where $m$ is any number less than $p-1$ and prime to it. Thus in the foregoing example $p=7$, where the prime roots were 3 and 5, the integers less than 6 and prime to it are 1, 5; and we, in fact, have $5\equiv3^5$ and $3\equiv5^5$ (mod.\ 7).

\textbf{18. Integers belonging to a given exponent; index of a number.} If, as before, $p-1=ef$, that is, if $f$ be a submultiple of $p-1$, then any integer $x$ such that $x^f$ is the lowest power of $x$ which is $\equiv1$ (mod.\ $p$) is said to belong to the exponent $f$. The number of residues, or terms of the series 1, 2, 3,\ldots, $p-1$, which belong to the exponent $f$ is $\phi(f)$, the number of integers less than $f$ and prime to it; these are the roots of the congruence $[x^f-1]=0$ of the order $\phi(f)$. It is hardly necessary to remark that the prime roots belong to the exponent $p-1$.

A number $x=g^m$ is said to have the index $m$; observe the distinction between the two terms exponent and index; and, further, that the index is dependent on the selected prime root $g$.

\textbf{19. Special forms of composite modulus.} If instead of a prime modulus $p$ we have a modulus $p^m$ which is the power of an odd prime, or a modulus $2p$ or $2p^m$ which is twice an odd prime or a power of an odd prime, then there is a theory analogous to that of prime roots, viz.\ the numbers less than the modulus and prime to it are congruent to successive powers of a prime root $g$; thus, if $p^m=3^2$, we have
\[
2,\ 4,\ 8,\ 16,\ 32,\ 64\equiv2,\ 4,\ 8,\ 7,\ 5,\ 1\pmod{9},
\]
and if $2p^m=2\cdot3^2$, we have
\[
5,\ 25,\ 125,\ 625,\ 3125,\ 15625\equiv5,\ 7,\ 11,\ 13,\ 17,\ 1\pmod{18}.
\]
As regards the even prime 2 and its powers: for the modulus 2 or 4 the theory of prime roots does not come into existence, and for the higher powers it is not applicable; thus with modulus $=8$ the numbers less than 8 and prime to it are 1, 3, 5, 7; and we have $3^2=5^2=7^2=1$ (mod.\ 8).

\textbf{20. Composite modulus} $N=a^\alpha b^\beta c^\gamma\ldots$ \textbf{no prime roots; irregularity.} In the general case of a composite modulus it has been seen that, if $x$ is any number less than $N$ and prime to it, then $x^{\phi(N)}-1\equiv0$ (mod.\ $N$). But, except in the above-mentioned cases $p^m$, $2p^m$, 2 or 4, there is not any number $a$ such that $a^{\phi(N)}$ is the first power of $a$ which is $\equiv1$; there is always some submultiple $l=\phi(N)/\sigma$ such that $a^l$ is the first

\end{document}
